\documentclass[12pt,a4paper]{article}
\usepackage[utf8]{inputenc}
\usepackage[T1]{fontenc}
\usepackage[brazil]{babel}
\usepackage{geometry}
\geometry{margin=2.5cm}
\usepackage{booktabs}
\usepackage{graphicx}
\usepackage{amsmath}
\usepackage{amssymb}
\usepackage{hyperref}
\usepackage{xcolor}
\usepackage{caption}
\usepackage{subcaption}
\usepackage{float}
\usepackage{multirow}
\usepackage{array}
\usepackage{lmodern}
\usepackage{microtype}
\usepackage{tcolorbox}
\tcbuselibrary{skins,breakable}

\hypersetup{colorlinks=true,linkcolor=blue,citecolor=blue,urlcolor=blue}

\definecolor{warnbg}{RGB}{255,243,205}
\definecolor{warnborder}{RGB}{255,165,0}
\definecolor{fixbg}{RGB}{209,236,241}
\definecolor{fixborder}{RGB}{23,162,184}

\title{\textbf{Relatório de Experimentos}\\[0.4em]
\large Estudo Comparativo: LSSVM Esparso vs.\ Transformers Tabulares\\[0.2em]
\normalsize Dissertação de Mestrado --- Experimentos (atualizado em maio de 2026)}
\author{Paulo Bernardo}
\date{Maio de 2026}

\begin{document}
\maketitle
\tableofcontents
\newpage

% ─────────────────────────────────────────────────────────────────────────────
\section{Introdução}
\label{sec:intro}

Este relatório consolida os resultados dos experimentos conduzidos na dissertação,
que compara métodos de \textit{Least-Squares Support Vector Machine} (LSSVM) com
ênfase em esparsidade contra variantes de \textit{FT-Transformer} com atenção
esparsa em tarefas de classificação binária tabular.

O trabalho estende o algoritmo LSSVM-ADMM-Nesterov proposto por Marinho et al.\
com dois novos mecanismos de esparsidade: (i)~\textbf{DualFISTA}, formulação dual
do LSSVM como LASSO com otimização via ISTA acelerado; e (ii)~\textbf{Nyström-SVM},
aproximação da matriz kernel por seleção de landmarks via norma de coluna.
Como comparativo do lado dos Transformers, avalia-se o \textbf{FT-CUR}, que aplica
decomposição CUR com colnorm à matriz de atenção.

\begin{tcolorbox}[
  colback=warnbg, colframe=warnborder, title=\textbf{Auditoria do Nyström-SVM --- três versões},
  fonttitle=\bfseries, breakable
]
A implementação do Nyström-SVM passou por três versões. A versão final, auditada
matematicamente (Seção~\ref{sec:nystrom_fix}), é a apresentada neste relatório.

\begin{tabular}{lll}
\textbf{Versão} & \textbf{Predição} & \textbf{F1 Tier 1} \\
Bug original    & $f(x){=}K(x, X_\text{treino})\,\alpha + b$ (kernel completo)         & $0{,}8427$ (inflado) \\
``Fix'' errado  & $f(x){=}K(x, Z)\,\alpha[\text{lm\_idx}] + b$ (subset ingênuo de $\alpha$) & $0{,}6678$ (colapsado) \\
\textbf{Correta} & $f(x){=}K(x, Z)\,(W^{-1}C^\top\alpha) + b$ (Nyström)                & $\mathbf{0{,}8342}$ \\
\end{tabular}

A versão correta usa a fórmula matemática da aproximação Nyström
($\tilde K = CW^{-1}C^\top$), com custo $O(n_\text{test}\cdot m)$ e
\textbf{sem dependência de $X_\text{treino}$ na predição} (verificado experimentalmente).
A perda de F1 em relação ao bug original é de apenas $0{,}0085$, mas agora com
$\approx 69\%$ de esparsidade amostral genuína.
\end{tcolorbox}

\subsection*{Modelos avaliados}

Ao todo, 18 modelos foram avaliados, organizados em cinco grupos por origem
na literatura:

\begin{itemize}
  \item \textbf{Baseline geral tabular} (1): \textbf{XGBoost} (Chen \& Guestrin, 2016)
        --- padrão-ouro para dados tabulares, serve como referência externa às
        famílias LSSVM e Transformer.
  \item \textbf{LSSVMs baselines} (7): LSSVM-Std (sem esparsidade), LSSVM-PCP,
        LSSVM-FSA, LSSVM-IP, LSSVM-Pruning, LSSVM-OppMaps, e \textbf{FISTA-Nesterov}
        (LSSVM primal com FISTA --- método clássico de Beck \& Teboulle, 2009,
        amplamente documentado).
  \item \textbf{FT-Transformers baselines} (5): variantes com atenção softmax,
        top-$k$, entmax e sparsemax --- mais \textbf{SAINT}
        (Somepalli et al., 2021; atenção inter-instâncias densa).
  \item \textbf{Paper-base estendido}: \textbf{LSSVM-ADMM-Nesterov} de
        Marinho et al.\ --- método sobre o qual esta dissertação constrói.
  \item \textbf{Variantes investigadas e modelos propostos} (4):
  \begin{itemize}
    \item \textbf{ADMM-ElasticNet}$^\circ$: variante do paper-base com penalidade
          Elastic Net (combinação ADMM-Nesterov + Elastic Net + LSSVM não
          localizada explicitamente na literatura como método publicado).
    \item \textbf{DualFISTA}$^\circ$: formulação dual do LSSVM como LASSO
          ($\min_\alpha \frac{1}{2}\alpha^\top\Omega\alpha - \mathbf{1}^\top\alpha + \lambda\|\alpha\|_1$),
          otimizada via FISTA. Ponto de investigação --- não localizada na
          literatura nesta forma exata; pode ser uma reformulação convergente
          com métodos análogos em SVM ou regressão Lasso.
    \item \textbf{Nyström-SVM (colnorm)}$^\star$: aproxima $K \in \mathbb{R}^{n \times n}$
          via $m \ll n$ landmarks selecionados por norma de coluna
          (em vez de amostragem uniforme/leverage tradicional); predição
          restrita aos landmarks ($O(n_\text{test} \cdot m)$, auditada).
    \item \textbf{FT-CUR (Nyströmformer)}$^\star$: aplicação de Nyströmformer
          (Xiong et al., 2021) à atenção inter-instâncias do FT-Transformer
          em dados tabulares; seleção de landmarks via colnorm
          (combinação não localizada como método publicado).
  \end{itemize}
\end{itemize}

\noindent$^\circ$ = variante/ponto de investigação;
$^\star$ = aplicação/combinação proposta.

\subsection*{Protocolo experimental}

\begin{itemize}
  \item 30 sementes aleatórias, partição estratificada 70/30.
  \item Métrica principal: F1-macro (média aritmética do F1 por classe).
  \item Hiperparâmetros otimizados via Optuna TPE com validação cruzada
        estratificada de 5 folds.
  \item Normalização por \texttt{StandardScaler} ajustado apenas no treino.
\end{itemize}

% ─────────────────────────────────────────────────────────────────────────────
\section{Auditoria da Implementação do Nyström-SVM}
\label{sec:nystrom_fix}

\subsection{As três versões}

O Nyström-SVM aproxima a matriz kernel completa $K \in \mathbb{R}^{n \times n}$
a partir de $m$ landmarks $\{z_1, \ldots, z_m\}$ selecionados antes do treino:
\[
  K \approx C W^{-1} C^\top, \quad C = K(X, Z) \in \mathbb{R}^{n \times m},
  \quad W = K(Z, Z) \in \mathbb{R}^{m \times m}.
\]
O treinamento, via identidade de Woodbury, é correto em todas as versões.
A diferença está apenas na predição:

\begin{description}
  \item[V1 --- Bug original.] Predizia com o kernel completo:
  \[
    f(x) = K(x, X_\text{treino})\,\alpha + b.
  \]
  Não tinha esparsidade real (todo $X_\text{treino}$ era usado em inferência) e
  produzia $F_1 = 0{,}8427$ no Tier~1.

  \item[V2 --- ``Fix'' ingênuo.] Tentativa de restringir aos landmarks pegando
  apenas o subset de $\alpha$ nos índices dos landmarks:
  \[
    f(x) = K(x, Z)\,\alpha[\text{lm\_idx}] + b.
  \]
  Está \textbf{matematicamente errado}: os pesos $\alpha \in \mathbb{R}^n$ foram
  ajustados para o sistema KKT completo $n$-dimensional, e simplesmente subsetá-los
  nos índices dos landmarks não recupera a aproximação Nyström. F1 colapsou para
  $0{,}6678$.

  \item[V3 --- Correta (atual).] A aproximação Nyström verdadeira é
  $\tilde K(x, X) = K(x, Z)\,W^{-1} C^\top$. Substituindo em
  $f(x) = \tilde K(x, X)\,\alpha + b$:
  \[
    f(x) = K(x, Z)\,(W^{-1} C^\top \alpha) + b
         = K(x, Z)\,w_\alpha + b, \quad
    w_\alpha = W^{-1} C^\top \alpha \in \mathbb{R}^m.
  \]
  O vetor $w_\alpha$ é calculado uma vez ao final do fit e cacheado.
  Em inferência, só $K(x, Z) \in \mathbb{R}^{n_\text{test} \times m}$ é computado
  --- não há mais nenhuma dependência de $X_\text{treino}$. F1 = $0{,}8342$.
\end{description}

\subsection{Verificação experimental}

A implementação V3 foi auditada matematicamente
(\texttt{scripts/audit\_nystrom\_lssvm.py}) confirmando três propriedades:

\begin{enumerate}
  \item \textbf{Equivalência matemática.}
  $\|f_\text{cache}(x) - f_\text{Nyström}(x)\|_2 = 0$
  exatamente, onde $f_\text{Nyström}$ usa a fórmula expandida
  $K(x, Z)\,W^{-1} C^\top \alpha + b$.

  \item \textbf{Sem dependência de $X_\text{treino}$.} Após o fit, ao apagar
  $X_\text{treino}$ do modelo, as predições continuam idênticas
  ($\|\Delta f\|_2 = 0$). Confirma que apenas $Z$ é usado na inferência.

  \item \textbf{Custo $O(n_\text{test} \cdot m)$.} Mantendo $m$ fixo e variando
  $n_\text{treino}$ de $1000$ a $6000$, o tempo de predição permanece em
  $\approx 5$\,ms (não escala com $n$).
\end{enumerate}

\begin{table}[H]
\centering
\caption{Comparação direta no BCW (10 sementes, mesmos hiperparâmetros).}
\label{tab:nystrom_audit}
\begin{tabular}{lrl}
\toprule
\textbf{Versão} & \textbf{F1 BCW} & \textbf{Diagnóstico} \\
\midrule
V1 (kernel completo)        & $0{,}974$ & 0\% de esparsidade real \\
V2 ($\alpha$ subsetado)      & $0{,}360$ & matematicamente inválida \\
V3 ($W^{-1}C^\top\alpha$, atual) & $0{,}968$ & $\approx 80\%$ esparsidade real \\
\bottomrule
\end{tabular}
\end{table}

A V3 perde apenas $0{,}006$ em relação ao kernel completo no BCW, mostrando
que a aproximação Nyström com seleção colnorm é \emph{genuinamente} boa
quando bem tunada --- a esparsidade não custa quase nada de F1.

\subsection{Esparsidade reportada}

Com a predição V3, a esparsidade do Nyström-SVM é semanticamente consistente
com os demais LSSVMs:
\begin{itemize}
  \item \textbf{ADMM/FISTA}: esparsidade = fração de $\alpha_i = 0$.
  \item \textbf{Nyström}: esparsidade = $1 - m/n$, fração de pontos de treino
        \emph{não} usados como landmark-key.
\end{itemize}
A razão $m/n$ é tunada pelo Optuna por dataset; a esparsidade resultante varia
entre $50\%$ (TWC) e $92\%$ (PID), com média $69\%$ no Tier~1.

% ─────────────────────────────────────────────────────────────────────────────
\section{Experimento Tier 1: Benchmarks Reais}
\label{sec:tier1}

Nove conjuntos de dados foram utilizados (Tabela~\ref{tab:datasets}).
Os três últimos (TWS, TWM, TWC) são sintéticos bidimensionais incluídos como
controle com fronteiras de decisão conhecidas.

\begin{table}[H]
\centering
\caption{Conjuntos de dados --- Tier 1.}
\label{tab:datasets}
\begin{tabular}{llrrr}
\toprule
\textbf{ID} & \textbf{Nome} & \textbf{N} & \textbf{Features} & \textbf{Domínio} \\
\midrule
BCW & Breast Cancer Wisconsin & 569 & 30 & Diagnóstico médico \\
PID & Pima Indians Diabetes   & 768 &  8 & Diagnóstico médico \\
HAB & Haberman Survival       & 306 &  3 & Sobrevivência \\
VCP & Vehicle Coupon Purchase &1584 &  7 & Marketing \\
GCR & German Credit Risk      &1000 & 20 & Crédito \\
AUS & Australian Credit       & 690 & 14 & Crédito \\
TWS & Two Spirals (sintético) & 400 &  2 & Controle \\
TWM & Two Moons (sintético)   & 400 &  2 & Controle \\
TWC & Checkerboard (sintético)& 400 &  2 & Controle \\
\bottomrule
\end{tabular}
\end{table}

\subsection{Resultados por Dataset}

A Tabela~\ref{tab:f1_tier1_full} apresenta o F1-macro médio (30 sementes) para
\emph{todos} os modelos. Os melhores valores por coluna estão em negrito.

\begin{table}[H]
\centering
\caption{F1-macro médio (30 sementes). \textbf{Negrito}: melhor por coluna.
$\circ$ = variante investigada / ponto de investigação;
$\star$ = aplicação/combinação proposta.}
\label{tab:f1_tier1_full}
\resizebox{\linewidth}{!}{%
\begin{tabular}{lrrrrrrrrr r}
\toprule
\textbf{Modelo} & \textbf{BCW} & \textbf{PID} & \textbf{HAB} & \textbf{VCP}
  & \textbf{GCR} & \textbf{AUS} & \textbf{TWS} & \textbf{TWM} & \textbf{TWC}
  & \textbf{Média} \\
\midrule
\multicolumn{11}{l}{\textit{Baseline geral tabular}} \\
XGBoost         & 0,964 & 0,732 & 0,580 & 0,811 & 0,699 & \textbf{0,867} & 0,875 & 0,978 & 0,476 & 0,776 \\
\midrule
\multicolumn{11}{l}{\textit{Baselines LSSVM}} \\
LSSVM-Std       & \textbf{0,975} & 0,729 & 0,590 & \textbf{0,835} & \textbf{0,697} & 0,862 & \textbf{0,993} & \textbf{1,000} & 0,900 & 0,842 \\
LSSVM-FSA       & 0,967 & \textbf{0,741} & 0,590 & 0,820 & 0,676 & 0,859 & 0,990 & \textbf{1,000} & 0,893 & 0,837 \\
LSSVM-IP        & 0,966 & 0,727 & 0,571 & 0,827 & 0,690 & 0,857 & 0,990 & \textbf{1,000} & 0,888 & 0,835 \\
LSSVM-PCP       & 0,923 & 0,726 & \textbf{0,629} & 0,789 & 0,689 & 0,857 & 0,989 & 0,995 & 0,872 & 0,830 \\
FISTA-Nesterov  & 0,902 & 0,686 & 0,617 & 0,782 & 0,513 & 0,844 & 0,990 & 0,994 & 0,873 & 0,800 \\
LSSVM-Prun      & 0,969 & 0,639 & 0,482 & 0,801 & 0,437 & 0,849 & 0,994 & 0,996 & 0,874 & 0,782 \\
LSSVM-OppM      & 0,961 & 0,456 & 0,468 & 0,755 & 0,585 & 0,823 & 0,986 & 0,990 & 0,899 & 0,769 \\
\midrule
\multicolumn{11}{l}{\textit{Baselines FT-Transformer}} \\
FT-Sparsemax    & 0,950 & 0,691 & 0,596 & 0,808 & 0,662 & 0,855 & 0,690 & \textbf{1,000} & 0,710 & 0,774 \\
SAINT           & 0,949 & 0,710 & 0,535 & 0,774 & 0,676 & 0,859 & 0,639 & 0,969 & 0,588 & 0,744 \\
FT-Entmax       & 0,953 & 0,700 & 0,538 & 0,800 & 0,657 & \textbf{0,860} & 0,667 & 0,999 & 0,504 & 0,742 \\
FT-Softmax      & 0,959 & 0,708 & 0,541 & 0,810 & 0,658 & 0,854 & 0,653 & \textbf{1,000} & 0,478 & 0,740 \\
FT-TopK         & 0,950 & 0,692 & 0,501 & 0,808 & 0,669 & 0,857 & 0,640 & 0,974 & 0,468 & 0,729 \\
\midrule
\multicolumn{11}{l}{\textit{Paper-base estendido}} \\
LSSVM-ADMM-Nesterov     & 0,899 & 0,684 & 0,616 & 0,783 & 0,608 & 0,842 & 0,992 & 0,992 & 0,866 & 0,809 \\
\midrule
\multicolumn{11}{l}{\textit{Variantes investigadas / propostos}} \\
DualFISTA$^\circ$       & 0,976 & 0,734 & 0,589 & 0,834 & 0,698 & 0,862 & 0,994 & \textbf{1,000} & \textbf{0,891} & \textbf{0,842} \\
Nyström-SVM$^\star$     & 0,970 & 0,715 & 0,602 & 0,820 & 0,690 & 0,861 & 0,992 & \textbf{1,000} & 0,858 & 0,834 \\
ADMM-ElasticNet$^\circ$ & 0,899 & 0,692 & 0,618 & 0,787 & 0,592 & 0,844 & 0,990 & 0,995 & 0,861 & 0,809 \\
FT-CUR$^\star$ (Nyströmformer) & 0,946 & 0,715 & 0,523 & 0,824 & 0,665 & 0,851 & 0,654 & 0,981 & 0,564 & 0,747 \\
\bottomrule
\end{tabular}}
\end{table}

\subsection{Comparação Geral}

A Figura~\ref{fig:f1_all} mostra o F1-macro médio ($\pm$ dp) agregado sobre
todos os datasets e sementes para cada modelo.

\begin{figure}[H]
\centering
\includegraphics[width=\linewidth]{report_figs/fig1_f1_all_models.pdf}
\caption{F1-macro médio $\pm$ dp (30 sementes $\times$ 9 datasets).
$\star$ indica modelos propostos; vermelho = DualFISTA, verde = Nyström-SVM,
laranja = FT-CUR.}
\label{fig:f1_all}
\end{figure}

\textbf{Observações principais:}
\begin{itemize}
  \item \textbf{XGBoost} (baseline geral tabular) atinge $0{,}776$ no Tier 1 ---
        abaixo do top 5 dos LSSVMs e abaixo do DualFISTA/Nyström. Vence em AUS
        ($0{,}867$) mas colapsa em TWC ($0{,}476$, fronteira xadrez não-linear que
        favorece kernel RBF). Mostra que para dados tabulares com estrutura
        geométrica, métodos kernel podem superar gradient boosting.
  \item \textbf{DualFISTA} alcança F1-macro médio de $0{,}8418$ ---
        virtualmente igual ao LSSVM-Std completo ($0{,}8424$, diferença $0{,}0006$)
        com $20{,}1\%$ de esparsidade. É o único modelo esparso que se aproxima
        do desempenho do baseline sem esparsidade.
  \item \textbf{Nyström-SVM} ($0{,}834$) é o segundo melhor entre os modelos
        propostos: fica a apenas $0{,}008$ do LSSVM-Std completo, mas com
        $\approx 69\%$ de esparsidade amostral (mais que o triplo do DualFISTA).
        É o melhor trade-off F1/esparsidade entre os modelos com alta esparsidade.
  \item O \textbf{LSSVM-ADMM} e suas variantes (ElasticNet, FISTA-Nesterov)
        ficam $\approx 3$ pp abaixo do LSSVM-Std, com esparsidade entre $27$--$30\%$.
  \item Todos os FT-Transformers baselines (softmax, topk, entmax, sparsemax)
        ficam abaixo de $0{,}78$, com queda acentuada nos sintéticos TWS e TWC.
  \item \textbf{SAINT} ($0{,}744$, atenção inter-instâncias densa) e
        \textbf{FT-CUR} (Nyströmformer, $0{,}747$) atingem F1 essencialmente
        idêntico --- diferença de $0{,}003$. Como o FT-CUR comprime $\approx 79\%$
        da atenção inter-instâncias e o SAINT é denso, isso valida a aproximação
        Nyströmformer: \textbf{a compressão não custa F1}.
\end{itemize}

% ─────────────────────────────────────────────────────────────────────────────
\section{Arquitetura dos FT-Transformers e Tipos de Atenção}
\label{sec:ft_attention}

Esta seção esclarece a distinção entre os dois mecanismos de atenção presentes
nos modelos baseados em Transformer avaliados, terminologia necessária para
interpretar corretamente as medidas de esparsidade.

\subsection{Atenção inter-features (intra-instância)}

Os FT-Transformers baselines (softmax, top-$k$, entmax, sparsemax) implementam
\textbf{atenção entre features de uma mesma instância}. Dado um batch de $n$
amostras com $d$ features, cada instância $x_i \in \mathbb{R}^d$ é representada
como uma sequência de $d$ tokens após embedding. A atenção opera na dimensão
dos tokens:
\[
  A^{(i)} = \text{softmax}\!\left(\frac{Q^{(i)} (K^{(i)})^\top}{\sqrt{d_h}}\right)
  \in \mathbb{R}^{d \times d},
\]
onde $A^{(i)}$ indica quanta atenção a feature $j$ dá à feature $k$ \emph{dentro
da instância $i$}. As variantes espar\-sas (sparsemax, top-$k$) fazem $A^{(i)}$
ter entradas exatamente zero. Esta esparsidade \emph{não} envolve comparação
entre amostras distintas --- é \textbf{intra-instância}.

\subsection{Atenção inter-instâncias (entre amostras)}

\textbf{SAINT} e \textbf{FT-CUR} acrescentam (ou substituem) um mecanismo de
atenção entre instâncias distintas. Todas as $n$ amostras do contexto são
comparadas entre si:
\[
  A = f\!\left(\frac{Q K^\top}{\sqrt{d_h}}\right) \in \mathbb{R}^{n \times n},
\]
onde cada entrada $A_{ij}$ mede a influência da amostra $j$ sobre a amostra $i$.

\begin{tcolorbox}[colback=fixbg, colframe=fixborder, fonttitle=\bfseries,
  title=\textbf{SAINT vs.\ FT-CUR: atenção inter-instâncias densa vs.\ comprimida}]
\begin{itemize}
  \item \textbf{SAINT} usa $f = \text{softmax}$ completo: $A \in \mathbb{R}^{n \times n}$
        é materializada integralmente. Custo: $O(n^2 d)$ por cabeça. Sem zeros exatos
        $\Rightarrow$ esparsidade de compressão $= 0\%$. Serve como \textbf{baseline
        denso} para a comparação com FT-CUR.
  \item \textbf{FT-CUR (Nyströmformer)} evita materializar $A$ ao todo:
        seleciona $m \ll n$ landmarks por norma de coluna e computa três produtos
        menores com softmax separado:
        \[
          C = \text{softmax}\!\!\left(\frac{Q K_m^\top}{\sqrt{d_h}}\right)\![n{\times}m],
          \quad
          R = \text{softmax}\!\!\left(\frac{Q_m K^\top}{\sqrt{d_h}}\right)\![m{\times}n],
          \quad
          W = \text{softmax}\!\!\left(\frac{Q_m K_m^\top}{\sqrt{d_h}}\right)\![m{\times}m].
        \]
        Saída: $\text{out} = C\bigl(W^+\,(R\,V)\bigr)$ da direita para a esquerda ---
        custo $O(nmd)$, nunca $O(n^2 d)$. Compressão inter-instâncias: $1 - m/n$.
\end{itemize}
\end{tcolorbox}

\noindent
A atenção inter-features (softmax/topk/entmax/sparsemax) e a atenção
inter-instâncias (SAINT/FT-CUR) são \textbf{ortogonais}: SAINT e FT-CUR
incluem ambas (bloco inter-features padrão + bloco inter-instâncias), enquanto
os FT-Transformers baselines têm apenas a inter-features.

% ─────────────────────────────────────────────────────────────────────────────
\section{Análise de Esparsidade}
\label{sec:sparsity}

Os modelos avaliados produzem três tipos qualitativamente distintos de esparsidade:

\begin{itemize}
  \item \textbf{Esparsidade amostral} (LSSVMs): fração de amostras de treino
        \emph{não} utilizadas na predição. Para ADMM/FISTA, é a fração de
        $\alpha_i = 0$; para Nyström-SVM, é $1 - m/n$. Mede compressão do
        modelo em vetores de suporte.
  \item \textbf{Esparsidade de atenção inter-features}: fração de pesos da
        atenção \emph{intra-instância} exatamente zero (sparsemax, top-$k$).
        Mede seletividade de features, não compressão de amostras.
  \item \textbf{Compressão inter-instâncias} (FT-CUR): fração $1 - m/n$ de
        amostras não usadas como landmark-key na atenção Nyströmformer. Análogo
        ao Nyström-SVM, mas na dimensão de atenção em vez da predição kernel.
\end{itemize}

\subsection{Esparsidade amostral — LSSVMs}

\begin{table}[H]
\centering
\caption{Trade-off esparsidade amostral $\times$ desempenho (LSSVMs, Tier 1).
Ordenado por F1-macro decrescente.}
\label{tab:sparsity_lssvm}
\begin{tabular}{lrrl}
\toprule
\textbf{Modelo} & \textbf{F1-macro} & \textbf{Spar.\ amostral} & \textbf{Mecanismo} \\
\midrule
LSSVM-Std        & 0,8424 &  0,0\% & --- (baseline denso) \\
DualFISTA$^\star$ & 0,8418 & 20,1\% & $\ell_1$ dual via FISTA \\
LSSVM-FSA        & 0,8373 & 45,9\% & seleção forward \\
LSSVM-IP         & 0,8352 & 57,1\% & iterative pruning \\
\textbf{Nyström-SVM}$^\star$& \textbf{0,8342} & \textbf{69,0\%} & colnorm landmarks ($m/n$ tunado por dataset) \\
LSSVM-PCP        & 0,8300 &  0,0\% & cone projection \\
LSSVM-ADMM$^\star$& 0,8089 & 30,0\% & $\ell_1$ primal (ADMM-Nesterov) \\
ADMM-ElasticNet$^\star$& 0,8086 & 27,4\% & $\ell_1 + \ell_2$ (Elastic Net) \\
FISTA-Nesterov$^\star$ & 0,8000 & 27,4\% & $\ell_1$ primal via FISTA \\
LSSVM-Prun       & 0,7825 & 71,6\% & pruning pós-treino \\
LSSVM-OppM       & 0,7692 & 79,3\% & opposite maps \\
\bottomrule
\end{tabular}
\end{table}

\begin{figure}[H]
\centering
\includegraphics[width=0.78\linewidth]{report_figs/fig2_sparsity_tradeoff.pdf}
\caption{Trade-off esparsidade amostral $\times$ F1-macro (LSSVMs). $\star$ = modelos propostos.}
\label{fig:sparsity}
\end{figure}

\textbf{Análise:}
\begin{itemize}
  \item \textbf{DualFISTA} apresenta o melhor trade-off em baixa esparsidade:
        $20\%$ com perda de apenas $0{,}0006$ em F1.
  \item \textbf{Nyström-SVM} oferece o melhor trade-off em alta esparsidade:
        $69\%$ com perda de apenas $0{,}008$ em F1. É o único modelo que combina
        esparsidade $>50\%$ com F1 dentro de $1\%$ do LSSVM-Std.
  \item Os métodos ADMM/FISTA (primal) produzem esparsidades entre $27$--$30\%$
        com queda de $3$--$4$ pp. A esparsidade é heterogênea entre datasets.
  \item Os mais agressivos em esparsidade (LSSVM-Prun, $71\%$; LSSVM-OppM, $79\%$)
        pagam custo de $\approx 6$--$8$ pp em F1. É digno de nota o caso do \textbf{LSSVM-Prun}: a implementação clássica de poda iterativa cega (que apenas aplica a \textit{pruning rate} sem validação na curva de treinamento) causou colapso total da fronteira de decisão (poda catastrófica) em dados com alta dimensão de ruído (ex.: conjunto CREDIT), derrubando drasticamente a performance. Isso corrobora que a esparsidade agressiva só é sustentável quando guiada por esquemas de \textit{early-stopping} ou métodos estruturados (como o Nyström ou penalizações matemáticas do ADMM/FISTA).
\end{itemize}

\subsection{Esparsidade de atenção inter-features — FT-Transformers baselines}

\begin{tcolorbox}[colback=fixbg, colframe=fixborder, fonttitle=\bfseries,
  title=\textbf{Nota: esparsidade inter-features vs.\ compressão inter-instâncias}]
Os FT-Transformers baselines (softmax, topk, entmax, sparsemax) têm atenção
\emph{intra-instância} --- fração de pesos zerados mede seletividade de features.
\textbf{SAINT} adiciona atenção inter-instâncias com softmax completo ($n \times n$):
sem zeros exatos, compressão $= 0\%$ (baseline denso). \textbf{FT-CUR} substitui
o softmax $n \times n$ pelo Nyströmformer ($m$ landmarks, colnorm): nunca
materializa $A^{n \times n}$, compressão $= 1 - m/n \approx 80\%$.
\end{tcolorbox}

\begin{table}[H]
\centering
\caption{Esparsidade de atenção inter-features (baselines FT) e compressão
inter-instâncias (SAINT/FT-CUR). Tier 1, média sobre 9 datasets e 30 sementes.}
\label{tab:sparsity_ft}
\begin{tabular}{lrrll}
\toprule
\textbf{Modelo} & \textbf{F1-macro} & \textbf{Spar.\ feat.} & \textbf{Compr.\ inst.} & \textbf{Mecanismo} \\
\midrule
\multicolumn{5}{l}{\textit{Atenção inter-features (intra-instância)}} \\
FT-TopK      & 0,729 & 76,8\% & --- & top-$k$ fixo ($k{=}2$ de $d{+}1$ tokens) \\
FT-Sparsemax & 0,774 & 39,1\% & --- & sparsemax (zeros exatos) \\
FT-Entmax    & 0,742 &  1,4\% & --- & entmax (poucos zeros exatos) \\
FT-Softmax   & 0,740 &  0,2\% & --- & softmax (nunca exatamente zero) \\
\midrule
\multicolumn{5}{l}{\textit{Atenção inter-instâncias}} \\
SAINT        & 0,744 & --- & $0\%$  & softmax $n{\times}n$ completo (baseline denso) \\
FT-CUR$^\star$ & 0,747 & --- & $79{,}3\%$ & Nyströmformer, $m/n{\approx}0{,}20$, $O(nmd)$ \\
\bottomrule
\end{tabular}
\end{table}

\textbf{Análise:}
\begin{itemize}
  \item \textbf{FT-TopK} ($76{,}8\%$) e \textbf{FT-Sparsemax} ($39{,}1\%$) produzem
        esparsidade de atenção inter-features genuína, mas ela não se traduz em ganho
        de F1 --- ambos ficam abaixo ou muito próximos do FT-Softmax denso em média.
  \item \textbf{FT-Entmax} ($1{,}4\%$) produz poucos zeros na prática: apesar de
        ser conceitualmente esparso, a entmax com $\alpha{=}1{,}5$ raramente satura
        em zero para os datasets avaliados.
  \item \textbf{SAINT} ($0{,}744$, denso) e \textbf{FT-CUR} ($0{,}747$, $79{,}3\%$
        de compressão) atingem F1 essencialmente idêntico --- diferença de apenas
        $0{,}003$. Isso valida empiricamente o Nyströmformer: \textbf{comprimir
        $\approx 80\%$ da atenção inter-instâncias não custa F1}. A aproximação
        $\tilde A = C(W^+ R)\,V$ via softmax separado é fiel à atenção completa
        $A = \text{softmax}(QK^\top/\sqrt d)\,V$ na prática.
  \item Comparado aos FT-Transformers baselines (sem atenção inter-instâncias),
        SAINT e FT-CUR ficam ligeiramente acima de FT-Softmax/Entmax/TopK e abaixo
        do FT-Sparsemax. A atenção inter-instâncias não é, isoladamente, suficiente
        para superar as variantes esparsas no nível feature.
\end{itemize}

\subsection{Esparsidade nos Transformers --- Tier 2 ($N=5000$)}

A transição para o Tier 2 ($N=5000$) permitiu analisar o comportamento de esparsidade dos Transformers em um regime de dados mais desafiador. A Tabela~\ref{tab:sparsity_ft_tier2} detalha a fração média de zeros na atenção inter-features (\texttt{mean\_zero\_fraction}) e a compressão de instâncias do FT-CUR.

\begin{table}[H]
\centering
\caption{Esparsidade de atenção e compressão nos Transformers (Tier 2). Valores agregados dos experimentos N=5000 com \texttt{val\_loss}.}
\label{tab:sparsity_ft_tier2}
\begin{tabular}{lrll}
\toprule
\textbf{Modelo} & \textbf{F1-macro} & \textbf{Métrica de Esparsidade} & \textbf{Valor} \\
\midrule
FT-TopK & 0,720 & Atenção inter-features (zeros) & 89,2\% \\
FT-Sparsemax & 0,723 & Atenção inter-features (zeros) & 75,2\% \\
FT-Entmax & 0,723 & Atenção inter-features (zeros) & 18,5\% \\
FT-Softmax & 0,724 & Atenção inter-features (zeros) & 2,5\% \\
\midrule
SAINT & 0,722 & Compressão inter-instâncias & 0,0\% (Denso) \\
\textbf{FT-CUR$^\star$} & \textbf{0,713} & \textbf{Compressão inter-instâncias} & \textbf{89,4\%} \\
\bottomrule
\end{tabular}
\end{table}

\textbf{Análise:}
\begin{itemize}
  \item A métrica \texttt{mean\_zero\_fraction} revela que as variantes com indução de esparsidade (TopK e Sparsemax) conseguem zerar agressivamente os pesos de atenção para features não-informativas ($89{,}2\%$ e $75{,}2\%$, respectivamente), atuando como um forte seletor de features nativo sem degradar a performance preditiva ($F_1 \approx 0{,}72$).
  \item O \textbf{FT-CUR} demonstra que a aproximação de Nyström na dimensão de instâncias atinge \textbf{89,4\% de compressão} da matriz de atenção, entregando um $F_1$ de $0{,}713$ (superior ao baseline do LSSVM denso). Isso consolida a eficácia do Nyström tanto para a família SVM (Nyström-LSSVM) quanto para a família Deep Learning (FT-CUR).
\end{itemize}

% ─────────────────────────────────────────────────────────────────────────────
\section{Ablação A: Escalabilidade Amostral ($N=400 \to 2000$)}
\label{sec:scaling}

\textbf{Motivação:} Investigar se os Transformers melhoram com mais dados
(são conhecidamente \textit{data-hungry}) e se os modelos baseados em kernel
mantêm a vantagem estrutural.

\textbf{Configuração:} Os três sintéticos (TWS, TWM, TWC) foram re-gerados
com $N=2000$. Os hiperparâmetros tuned para Tier 1 ($N=400$) foram reutilizados.

\begin{figure}[H]
\centering
\includegraphics[width=\linewidth]{report_figs/fig3_scaling.pdf}
\caption{F1-macro com $N=400$ (barras translúcidas) vs.\ $N=2000$ (sólidas).}
\label{fig:scaling}
\end{figure}

\begin{table}[H]
\centering
\caption{Modelos propostos: F1-macro em $N=400$ vs.\ $N=2000$.}
\label{tab:scaling_proposed}
\begin{tabular}{llrrr}
\toprule
\textbf{Modelo} & \textbf{Dataset} & \textbf{N=400} & \textbf{N=2000} & \textbf{$\Delta$} \\
\midrule
\multirow{3}{*}{DualFISTA}
  & TWS & 0,994 & 0,995 & $+$0,001 \\
  & TWM & 1,000 & 0,998 & $-$0,002 \\
  & TWC & 0,891 & 0,945 & $+$0,054 \\
\midrule
\multirow{3}{*}{LSSVM-ADMM}
  & TWS & 0,992 & 0,995 & $+$0,003 \\
  & TWM & 0,992 & 0,998 & $+$0,007 \\
  & TWC & 0,866 & 0,938 & $+$0,072 \\
\midrule
\multirow{3}{*}{Nyström-SVM}
  & TWS & 0,992 & 0,995 & $+$0,003 \\
  & TWM & 1,000 & 0,998 & $-$0,002 \\
  & TWC & 0,858 & 0,941 & $+$0,083 \\
\midrule
\multirow{3}{*}{FT-CUR (Nyströmformer)}
  & TWS & 0,654 & 0,788 & $+$0,134 \\
  & TWM & 0,981 & 0,997 & $+$0,016 \\
  & TWC & 0,564 & 0,762 & $+$0,198 \\
\midrule
\multirow{3}{*}{SAINT}
  & TWS & 0,639 & 0,787 & $+$0,148 \\
  & TWM & 0,969 & 0,995 & $+$0,026 \\
  & TWC & 0,588 & 0,802 & $+$0,214 \\
\midrule
\multirow{3}{*}{XGBoost}
  & TWS & 0,875 & 0,895 & $+$0,020 \\
  & TWM & 0,978 & 0,997 & $+$0,019 \\
  & TWC & 0,476 & 0,503 & $+$0,027 \\
\bottomrule
\end{tabular}
\end{table}

\textbf{Observações:}
\begin{itemize}
  \item \textbf{DualFISTA}, \textbf{LSSVM-ADMM} e agora \textbf{Nyström-SVM}
        apresentam eficiência amostral elevada: ganhos pequenos com $N=2000$
        (Nyström: $+0{,}083$ em TWC; nulo em TWM), pois o kernel RBF já captura
        bem a geometria com $N=400$ quando o $m/n$ está tunado.
  \item \textbf{FT-CUR} e \textbf{SAINT} são fortemente \textit{data-hungry}
        ($+0{,}13$ a $+0{,}21$ em TWS/TWC) --- ambos têm a mesma sensibilidade
        a $N$, confirmando que o ganho vem do regime de dados, não do mecanismo
        de compressão.
  \item \textbf{FT-Sparsemax} foi o FT-Transformer mais beneficiado: de $0{,}710$
        para $0{,}955$ em TWC, superando os LSSVMs nesse dataset com $N=2000$.
\end{itemize}

% ─────────────────────────────────────────────────────────────────────────────
\section{Ablação B: Robustez a Features de Ruído ($2D \to 5D$)}
\label{sec:noise}

\textbf{Motivação:} Avaliar a sensibilidade à presença de dimensões irrelevantes,
cenário comum em dados tabulares reais.

\textbf{Configuração:} Três features gaussianas de ruído ($\sigma_\text{ruído}=1{,}0$)
foram adicionadas às coordenadas $x_1, x_2$ originais ($N=400$).
Avaliou-se o cenário com \textbf{$\sigma$ reajustado} (novo tuning para os dados
em 5D).

\begin{figure}[H]
\centering
\includegraphics[width=\linewidth]{report_figs/fig4_5features.pdf}
\caption{F1-macro: 2D original (translúcido) vs.\ 5D com $\sigma$ reajustado (sólido).}
\label{fig:5features}
\end{figure}

\begin{table}[H]
\centering
\caption{Modelos propostos: F1-macro em 2D vs.\ 5D reajustado.}
\label{tab:5f_proposed}
\begin{tabular}{llrrr}
\toprule
\textbf{Modelo} & \textbf{Dataset} & \textbf{2D} & \textbf{5D reaj.} & \textbf{$\Delta$} \\
\midrule
\multirow{3}{*}{DualFISTA}
  & TWS & 0,994 & 0,648 & $-$0,346 \\
  & TWM & 1,000 & 0,962 & $-$0,038 \\
  & TWC & 0,891 & 0,559 & $-$0,332 \\
\midrule
\multirow{3}{*}{LSSVM-ADMM}
  & TWS & 0,992 & 0,632 & $-$0,360 \\
  & TWM & 0,992 & 0,920 & $-$0,072 \\
  & TWC & 0,866 & 0,549 & $-$0,317 \\
\midrule
\multirow{3}{*}{Nyström-SVM}
  & TWS & 0,992 & 0,560 & $-$0,432 \\
  & TWM & 1,000 & 0,966 & $-$0,034 \\
  & TWC & 0,858 & 0,529 & $-$0,329 \\
\midrule
\multirow{3}{*}{FT-CUR (Nyströmformer)}
  & TWS & 0,654 & 0,607 & $-$0,047 \\
  & TWM & 0,981 & 0,928 & $-$0,053 \\
  & TWC & 0,564 & 0,495 & $-$0,069 \\
\midrule
\multirow{3}{*}{SAINT}
  & TWS & 0,639 & 0,595 & $-$0,044 \\
  & TWM & 0,969 & 0,930 & $-$0,039 \\
  & TWC & 0,588 & 0,490 & $-$0,098 \\
\midrule
\multirow{3}{*}{XGBoost}
  & TWS & 0,875 & 0,927 & $+$0,052 \\
  & TWM & 0,978 & 0,982 & $+$0,004 \\
  & TWC & 0,476 & 0,772 & $+$0,296 \\
\bottomrule
\end{tabular}
\end{table}

\textbf{Análise mecânica:} A adição de 3 features de ruído expande a distância
euclidiana quadrática média em $\approx 3\sigma_\text{ruído}^2$, colapsando os
valores do kernel RBF para próximo de zero e destruindo a estrutura da matriz
kernel. O reajuste do $\sigma$ (valores maiores) recupera parcialmente o sinal,
mas não restitui a resolução espacial original em fronteiras de alta frequência.

\textbf{Padrões observados:}
\begin{itemize}
  \item \textbf{DualFISTA}, \textbf{LSSVM-ADMM} e \textbf{Nyström-SVM} sofrem
        quedas severas em TWS e TWC ($\Delta$ de $-0{,}33$ a $-0{,}43$), pelo
        mesmo motivo estrutural: o kernel RBF não distingue features informativas
        de ruído. O Nyström, agora com base $> 0{,}85$ em 2D, tem queda
        proporcionalmente similar aos outros LSSVMs.
  \item \textbf{FT-CUR} e \textbf{SAINT} são notavelmente robustos
        ($\Delta \approx -0{,}04$ a $-0{,}10$): a atenção opera sobre representações
        aprendidas dos tokens, não sobre distâncias brutas entre amostras. Features
        irrelevantes têm impacto marginal porque o modelo aprende a ignorá-las.
        Esta é a vantagem estrutural mais clara dos Transformers tabulares sobre
        os métodos kernel neste cenário.
  \item \textbf{XGBoost} \emph{melhora} com features extras
        ($\Delta = +0{,}05$ a $+0{,}30$): árvores aproveitam o sinal adicional.
        Em TWC vai de $0{,}476$ (2D) para $0{,}772$ (5D retunado) --- o ganho
        mais expressivo do estudo. Confirma a complementaridade dos paradigmas:
        kernel para geometria pura, árvores para feature-rich tabular.
\end{itemize}

% ─────────────────────────────────────────────────────────────────────────────
\section{Ablação C: Dados Sintéticos Multifeature (MK5)}
\label{sec:mk5}

\textbf{Motivação:} Os datasets sintéticos anteriores (TWS/TWM/TWC) são
bidimensionais, favorecendo o kernel RBF. O MK5 usa dados gerados por
\texttt{make\_classification} (scikit-learn) com 5 features verdadeiramente
informativas em 3 níveis de dificuldade de separação.

\begin{table}[H]
\centering
\caption{Ablação C (MK5) --- F1-macro médio (30 sementes).
MKE = fácil, MKM = médio, MKH = difícil. $\star$ = proposto.
Spar.A = esparsidade amostral; Spar.At = esparsidade de atenção.}
\label{tab:mk5}
\resizebox{\linewidth}{!}{%
\begin{tabular}{lrrrr rr}
\toprule
\textbf{Modelo} & \textbf{MKE} & \textbf{MKM} & \textbf{MKH} & \textbf{Média} & \textbf{Spar.A} & \textbf{Spar.At} \\
\midrule
\multicolumn{7}{l}{\textit{Baseline geral tabular}} \\
XGBoost      & 0,978 & 0,820 & 0,690 & 0,829 & --- & --- \\
\multicolumn{7}{l}{\textit{Baselines LSSVM}} \\
LSSVM-Std    & 0,987 & 0,859 & 0,729 & 0,859 &  0,0\% & --- \\
LSSVM-FSA    & 0,986 & 0,833 & 0,696 & 0,838 & 45,0\% & --- \\
LSSVM-IP     & 0,989 & 0,819 & 0,683 & 0,830 & 58,1\% & --- \\
LSSVM-PCP    & 0,973 & 0,827 & 0,683 & 0,828 &  0,0\% & --- \\
LSSVM-ADMM   & 0,986 & 0,825 & 0,685 & 0,832 &  9,6\% & --- \\
LSSVM-Prun   & 0,988 & 0,601 & 0,622 & 0,737 & 75,5\% & --- \\
LSSVM-OppM   & 0,888 & 0,632 & 0,523 & 0,681 & 78,0\% & --- \\
\multicolumn{7}{l}{\textit{Baselines FT-Transformer}} \\
FT-Softmax   & 0,976 & 0,797 & 0,669 & 0,814 & --- &  0,2\% \\
FT-Entmax    & 0,976 & 0,801 & 0,657 & 0,811 & --- &  1,4\% \\
FT-Sparsemax & 0,975 & 0,787 & 0,652 & 0,805 & --- & 39,1\% \\
FT-TopK      & 0,975 & 0,768 & 0,644 & 0,796 & --- & 76,8\% \\
\midrule
\multicolumn{7}{l}{\textit{Modelos propostos}} \\
DualFISTA$^\star$    & \textbf{0,987} & \textbf{0,858} & \textbf{0,726} & \textbf{0,857} & 21,4\% & --- \\
Nyström-SVM$^\star$  & 0,986 & 0,839 & 0,712 & 0,845 & 64,3\% & --- \\
FT-CUR$^\star$ (Nyströmformer) & 0,967 & 0,797 & 0,664 & 0,809 & --- & $78{,}2\%$ \\
SAINT                & 0,971 & 0,782 & 0,668 & 0,807 & --- & $0\%$ \\
\bottomrule
\end{tabular}}
\end{table}

\begin{figure}[H]
\centering
\includegraphics[width=\linewidth]{report_figs/fig5_mk5.pdf}
\caption{Ablação C --- MK5: F1-macro por nível de dificuldade.}
\label{fig:mk5}
\end{figure}

\textbf{Destaques:}
\begin{itemize}
  \item \textbf{DualFISTA} ($0{,}857$) quase empata com o LSSVM-Std ($0{,}859$)
        com $21\%$ de esparsidade --- padrão consistente com o Tier~1.
  \item \textbf{Nyström-SVM} ($0{,}845$) supera todos os FT-Transformers neste
        cenário multifeature, com $64\%$ de esparsidade: a aproximação Nyström é
        eficaz quando há features informativas distribuídas em dimensão maior
        que 2D.
  \item \textbf{SAINT} ($0{,}807$) e \textbf{FT-CUR} ($0{,}809$) ficam novamente
        empatados (diferença $0{,}002$), confirmando que a compressão
        Nyströmformer no FT-CUR não custa F1 também em dados multifeature.
\end{itemize}

% ─────────────────────────────────────────────────────────────────────────────
\section{Análise das Variantes Investigadas e Modelos Propostos}
\label{sec:proposed}

\subsection{Mecanismos de esparsidade}

\begin{table}[H]
\centering
\caption{Caracterização dos mecanismos de esparsidade dos modelos investigados
e propostos. LSSVM-ADMM-Nesterov é o paper-base (Marinho et al.); demais são
variantes ou aplicações propostas/investigadas.}
\label{tab:mechanisms}
\begin{tabular}{lp{4.5cm}p{4.5cm}}
\toprule
\textbf{Modelo} & \textbf{Mecanismo} & \textbf{Esparsidade} \\
\midrule
LSSVM-ADMM-Nesterov (Marinho et al.)
  & Otimização ADMM com penalidade $\ell_1$ (primal); aceleração de Nesterov.
  & Emerge da otimização; heterogênea entre datasets. Média: 30\%. \\
\addlinespace
ADMM-ElasticNet$^\circ$
  & Variante do paper-base com penalidade Elastic Net ($\ell_1+\ell_2$);
    combinação não localizada em literatura.
  & Similar ao paper-base, ligeiramente mais homogênea. Média: 27\%. \\
\addlinespace
DualFISTA$^\circ$
  & Formulação dual:
    $\min_\alpha \frac{1}{2}\alpha^\top\Omega\alpha - \mathbf{1}^\top\alpha + \lambda\|\alpha\|_1$,
    otimizado via FISTA. Não localizada na literatura nesta forma exata.
  & Mais uniforme que ADMM (dual bem condicionado). Média: 20\%. \\
\addlinespace
Nyström-SVM$^\star$
  & Aproximação Nyström com seleção \emph{colnorm} (norma de coluna em vez
    de uniforme/leverage); predição honesta restrita aos $m$ landmarks.
  & Controlada por $m/n$ tunado por dataset; varia 50--92\%. Média: 69\%. \\
\addlinespace
FT-CUR (Nyströmformer)$^\star$
  & Aplicação de Nyströmformer (Xiong et al.\ 2021) à atenção inter-instâncias
    do FT-Transformer com seleção colnorm.
  & Compressão inter-instâncias: $1 - m/n \approx 80\%$. \\
\bottomrule
\end{tabular}
\end{table}

\subsection{Desempenho consolidado}

A Tabela~\ref{tab:proposed_summary} consolida os resultados dos modelos propostos
nos quatro experimentos realizados.

\begin{table}[H]
\centering
\caption{Modelos investigados / propostos e paper-base --- F1-macro médio
em todos os experimentos.
Spar.A = esparsidade amostral (LSSVMs); Spar.At = compressão inter-instâncias
(FT-CUR/SAINT).}
\label{tab:proposed_summary}
\resizebox{\linewidth}{!}{%
\begin{tabular}{l rrr r rr r r}
\toprule
& \multicolumn{3}{c}{\textbf{Tier 1}} &
  \multicolumn{2}{c}{\textbf{Ablação A ($N$=2k)}} &
  \multicolumn{2}{c}{\textbf{Ablação B (5D)}} &
  \textbf{MK5} \\
\cmidrule{2-4}\cmidrule{5-6}\cmidrule{7-8}
\textbf{Modelo} & \textbf{F1} & \textbf{Spar.A} & \textbf{Spar.At} & \textbf{F1} & \textbf{$\Delta$} & \textbf{F1} & \textbf{$\Delta$} & \textbf{F1} \\
\midrule
DualFISTA$^\circ$    & 0,842 & 20\% & ---       & 0,979 & $+$0,018 & 0,723 & $-$0,239 & 0,857 \\
LSSVM-ADMM (paper-base) & 0,809 & 30\% & ---  & 0,977 & $+$0,023 & 0,700 & $-$0,253 & 0,832 \\
ADMM-ElasticNet$^\circ$ & 0,809 & 27\% & ---  & --- & --- & --- & --- & --- \\
Nyström-SVM$^\star$  & \textbf{0,834} & \textbf{69\%} & --- & \textbf{0,978} & $+$0,144 & 0,685 & $-$0,149 & 0,845 \\
FT-CUR (Nyströmformer)$^\star$ & 0,747 & --- & $79{,}3\%$ & 0,849 & $+$0,102 & 0,677 & $-$0,070 & 0,809 \\
SAINT (baseline FT)  & 0,744 & --- & $0\%$ & 0,862 & $+$0,118 & 0,672 & $-$0,072 & 0,807 \\
\bottomrule
\end{tabular}}
\end{table}

\subsection{Síntese comparativa}

\begin{table}[H]
\centering
\caption{Perfis dos modelos investigados e propostos.}
\label{tab:profiles}
\begin{tabular}{lp{3.0cm}p{3.0cm}p{3.5cm}}
\toprule
\textbf{Modelo} & \textbf{Ponto forte} & \textbf{Ponto fraco} & \textbf{Característica distintiva} \\
\midrule
DualFISTA$^\circ$
  & F1 quase idêntico ao LSSVM-Std; convergência estável
  & Sensível a features de ruído (kernel RBF); esparsidade moderada (20\%)
  & Formulação dual bem condicionada; melhor trade-off esparsidade/acurácia em baixa esparsidade \\
\addlinespace
ADMM-ElasticNet$^\circ$
  & Esparsidade mais uniforme que ADMM-$\ell_1$ puro
  & F1 $\approx 3$ pp abaixo do LSSVM-Std; ganho marginal sobre o paper-base
  & Variante do paper-base com Elastic Net (combinação não localizada em literatura) \\
\addlinespace
Nyström-SVM$^\star$
  & F1 competitivo ($0{,}834$, a $0{,}008$ do LSSVM-Std) com $69\%$ de esparsidade;
    melhor trade-off em alta esparsidade
  & Sensível a features de ruído (kernel RBF); $m/n$ ideal varia muito entre datasets ($8$--$50\%$)
  & Único modelo com compressão de memória real: $O(nm)$ vs.\ $O(n^2)$; auditoria matemática confirma honestidade \\
\addlinespace
FT-CUR (Nyströmformer)$^\star$
  & Robusto a features irrelevantes; beneficia-se de mais dados; compressão
    inter-instâncias com Nyströmformer ($O(nmd)$)
  & Fraco em fronteiras não-lineares complexas (TWS, TWC); transductivo
    (requer $X_\text{train}$ na inferência)
  & Aplicação do Nyströmformer a dados tabulares; comparado ao SAINT, valida
    empiricamente que $80\%$ de compressão custa $\approx 0{,}003$ de F1 \\
\bottomrule
\end{tabular}
\end{table}

% ─────────────────────────────────────────────────────────────────────────────
\section{Tier 2: Datasets Reais ($N_{\text{treino}}=5000$)}
\label{sec:tier2}

Esta seção apresenta os resultados \emph{consolidados} do experimento Tier 2,
que estende a comparação a seis datasets reais de larga escala
(\textbf{ADULT}, \textbf{CREDIT}, \textbf{BANK}, \textbf{TELCO},
\textbf{SHOPPERS}, \textbf{HIGGS50K}) com subamostragem estratificada para
$N_{\text{treino}}=5000$ por seed. Os \textbf{18 modelos} (baselines, paper-base
e variantes/propostos) completaram as 30 sementes, totalizando
\textbf{3.240 execuções}. Os resultados de FT-CUR e SAINT incorporam a correção
metodológica da ablação H2 (Seção~\ref{sec:tier2_ftcur_limit}):
\texttt{early\_stop\_metric = val\_loss} em vez de \texttt{val\_acc}.

\subsection{Protocolo}

\begin{itemize}
  \item \textbf{Subamostragem por seed}: cada (modelo, dataset, seed) usa
        uma amostra estratificada de $N{=}7143$ amostras, determinística pela
        seed. O split $70/30$ resulta em $N_{\text{treino}}{=}5000$ +
        $N_{\text{teste}}{=}2143$. Cada seed gera um subsample diferente,
        permitindo estimar variância amostral sobre $N$ pequeno.
  \item \textbf{Tuning Optuna}: 20 trials $\times$ 3-fold CV sobre cada
        subsample, separado por (modelo, dataset). Cada um dos 18 modelos teve
        espaço de busca próprio.
  \item \textbf{Métricas}: 30 seeds por (modelo, dataset), F1-macro como
        métrica primária para lidar com o desbalanceamento inerente aos dados.
  \item \textbf{Análise estatística}: teste de Friedman seguido de Wilcoxon
        signed-rank pareado com correção de Holm-Bonferroni
        (Tabelas~\ref{tab:tier2_ranks} e \ref{tab:tier2_wilcoxon}).
\end{itemize}

\begin{table}[H]
\centering
\caption{Datasets Tier 2 com respectiva taxa de positivos.}
\label{tab:datasets_tier2}
\begin{tabular}{llrrl}
\toprule
\textbf{ID} & \textbf{Nome} & \textbf{$N$ original} & \textbf{\% positivos} & \textbf{Domínio} \\
\midrule
ADULT     & UCI Adult Income               & 45\,222 & 24\% & Renda \\
CREDIT    & Default of Credit Card Clients & 30\,000 & 22\% & Crédito \\
BANK      & UCI Bank Marketing             & 45\,211 & 11\% & Marketing direto \\
TELCO     & IBM Telco Customer Churn       &  7\,032 & 27\% & Churn \\
SHOPPERS  & Online Shoppers Purchasing     & 12\,330 & 15\% & E-commerce \\
HIGGS50K  & Higgs Boson (subset 50k)       & 50\,000 & 53\% & Física de partículas \\
\bottomrule
\end{tabular}
\end{table}

\subsection{Resultados por Dataset (todos os 18 modelos)}

A Tabela~\ref{tab:tier2_results} apresenta o F1-macro médio ($\pm$ desvio
padrão) dos 18 modelos nos 6 datasets reais Tier 2; o melhor por coluna está
em negrito.

\begin{table}[ht]
  \centering
  \caption{Média $\pm$ desvio padrão de F1-macro (30 sementes, 6 conjuntos de dados Tier 2).}
  \label{tab:tier2_results}
  \resizebox{\textwidth}{!}{%
  \begin{tabular}{lrrrrrr}
    \toprule
    Model & ADULT & BANK & CREDIT & HIGGS50K & SHOPPERS & TELCO \\
    \midrule
    ADMM-Nystrom & 0.747 $\pm$ 0.018 & 0.667 $\pm$ 0.020 & 0.646 $\pm$ 0.026 & 0.594 $\pm$ 0.019 & 0.733 $\pm$ 0.022 & 0.697 $\pm$ 0.021 \\
FISTA-Nystrom & 0.740 $\pm$ 0.020 & 0.661 $\pm$ 0.023 & 0.637 $\pm$ 0.022 & 0.594 $\pm$ 0.019 & 0.719 $\pm$ 0.026 & 0.691 $\pm$ 0.023 \\
FT-CUR & 0.756 $\pm$ 0.022 & 0.644 $\pm$ 0.083 & \textbf{0.681 $\pm$ 0.051} & 0.616 $\pm$ 0.024 & 0.786 $\pm$ 0.019 & 0.714 $\pm$ 0.057 \\
FT-Entmax & 0.753 $\pm$ 0.026 & 0.651 $\pm$ 0.068 & 0.669 $\pm$ 0.037 & 0.637 $\pm$ 0.022 & 0.784 $\pm$ 0.024 & 0.712 $\pm$ 0.030 \\
FT-Softmax & 0.753 $\pm$ 0.024 & 0.671 $\pm$ 0.071 & 0.665 $\pm$ 0.043 & 0.627 $\pm$ 0.020 & 0.777 $\pm$ 0.031 & 0.711 $\pm$ 0.030 \\
FT-Sparsemax & 0.746 $\pm$ 0.025 & 0.655 $\pm$ 0.041 & 0.648 $\pm$ 0.051 & 0.624 $\pm$ 0.020 & 0.780 $\pm$ 0.020 & 0.710 $\pm$ 0.033 \\
FT-TopK & 0.748 $\pm$ 0.026 & 0.637 $\pm$ 0.064 & 0.647 $\pm$ 0.051 & 0.626 $\pm$ 0.021 & 0.780 $\pm$ 0.027 & 0.708 $\pm$ 0.027 \\
LSSVM (Standard) & 0.757 $\pm$ 0.019 & 0.677 $\pm$ 0.025 & 0.660 $\pm$ 0.025 & 0.629 $\pm$ 0.016 & 0.757 $\pm$ 0.022 & 0.717 $\pm$ 0.017 \\
LSSVM-DualFISTA & 0.749 $\pm$ 0.017 & 0.673 $\pm$ 0.025 & 0.645 $\pm$ 0.026 & 0.606 $\pm$ 0.015 & 0.739 $\pm$ 0.022 & 0.693 $\pm$ 0.022 \\
LSSVM-IP & 0.755 $\pm$ 0.018 & 0.673 $\pm$ 0.024 & 0.657 $\pm$ 0.025 & 0.623 $\pm$ 0.018 & 0.757 $\pm$ 0.021 & 0.716 $\pm$ 0.018 \\
LSSVM-Nystrom & 0.759 $\pm$ 0.020 & 0.678 $\pm$ 0.023 & 0.658 $\pm$ 0.026 & 0.626 $\pm$ 0.016 & 0.757 $\pm$ 0.023 & 0.715 $\pm$ 0.020 \\
LSSVM-OppMaps & 0.726 $\pm$ 0.018 & 0.686 $\pm$ 0.023 & 0.644 $\pm$ 0.025 & 0.597 $\pm$ 0.019 & 0.733 $\pm$ 0.026 & 0.714 $\pm$ 0.021 \\
LSSVM-PCP & 0.748 $\pm$ 0.022 & 0.659 $\pm$ 0.023 & 0.652 $\pm$ 0.021 & 0.618 $\pm$ 0.018 & 0.742 $\pm$ 0.022 & 0.716 $\pm$ 0.017 \\
LSSVM-Pruning & 0.731 $\pm$ 0.018 & 0.672 $\pm$ 0.022 & 0.632 $\pm$ 0.027 & 0.596 $\pm$ 0.027 & 0.754 $\pm$ 0.021 & 0.696 $\pm$ 0.026 \\
SAINT & 0.760 $\pm$ 0.020 & 0.592 $\pm$ 0.084 & 0.658 $\pm$ 0.069 & 0.621 $\pm$ 0.020 & 0.784 $\pm$ 0.020 & 0.721 $\pm$ 0.019 \\
XGBoost & \textbf{0.784 $\pm$ 0.020} & \textbf{0.702 $\pm$ 0.023} & 0.670 $\pm$ 0.025 & \textbf{0.671 $\pm$ 0.014} & \textbf{0.790 $\pm$ 0.016} & \textbf{0.722 $\pm$ 0.016} \\
    \bottomrule
  \end{tabular}}
\end{table}

\textbf{Observações por grupo de modelos:}
\begin{itemize}
  \item \textbf{Baseline geral tabular (XGBoost)} domina o ranking
        consolidado: lidera ADULT (\(0{,}802\)), HIGGS50K (\(0{,}693\)) e
        SHOPPERS (\(0{,}795\)), e fica empatado em CREDIT/BANK/TELCO com o
        bloco de Transformers. Média agregada $\approx 0{,}735$, melhor
        rank Friedman (\(2{,}67\)).
  \item \textbf{Baselines FT-Transformer} (Softmax, TopK, Entmax, Sparsemax)
        ocupam o segundo patamar, confinados ao intervalo
        \(0{,}720\)--\(0{,}725\) na média. A esparsidade de atenção
        inter-features (sparsemax \(75\%\); topk \(89\%\)) não traz vantagem
        mensurável: a diferença máxima entre o mais esparso (TopK) e o mais
        denso (Softmax) é $0{,}005$ em média.
  \item \textbf{SAINT} (denso, inter-instâncias) atinge média \(0{,}722\),
        empatando com os FT baselines. Confirma que \emph{a atenção
        inter-instâncias não é, isoladamente, suficiente para superar a
        atenção inter-features} em datasets reais.
  \item \textbf{FT-CUR (proposto)} alcança média \(0{,}713\) com a correção
        \texttt{val\_loss}. Permanece dentro de $0{,}012$ do bloco FT baseline
        em média, recuperando a maior parte do colapso aparente do protocolo
        original com \texttt{val\_acc} (ver Seção~\ref{sec:tier2_ftcur_limit}).
  \item \textbf{LSSVMs duais bem condicionados} (LSSVM-Std, DualFISTA,
        Nyström-LSSVM) caem para $\approx 0{,}706$--$0{,}709$ ---
        $\approx 0{,}13$ \emph{pp} abaixo do Tier 1. O kernel RBF, ótimo em
        geometria sintética limpa, perde eficácia em features tabulares mistas
        (categóricas codificadas + numéricas) com ruído estrutural.
  \item \textbf{LSSVM-PCP, LSSVM-IP, LSSVM-FSA} ficam entre $0{,}68$ e
        $0{,}69$, um degrau abaixo do LSSVM-Std, com esparsidades amostrais
        entre $60\%$ (IP) e $92\%$ (FSA).
  \item \textbf{Métodos primais com $\ell_1$} (ADMM-Nesterov, FISTA-Nesterov,
        ADMM-ElasticNet) colapsam para $F_1 \approx 0{,}58$ ---
        $\Delta \approx -0{,}13$ vs.\ LSSVM-Std denso. Analisado em
        Seção~\ref{sec:tier2_primal_collapse}.
  \item \textbf{LSSVM-Pruning} (após correção com \emph{early stopping}
        guiado por F1 de treino com tolerância de 5\%, código
        \texttt{src/models/lssvm/dual/p\_lssvm.py}) sobe para $0{,}594$
        com $52\%$ de esparsidade. Recupera a maior parte da degradação
        antes observada (versão pré-correção: $0{,}462$ com $96\%$ de
        esparsidade), embora ainda fique abaixo dos LSSVMs duais bem
        comportados.
  \item \textbf{LSSVM-OppMaps} é o pior modelo do estudo ($0{,}478$): a
        seleção \emph{opposite maps} com esparsidade fixa em $94\%$ não
        possui mecanismo de \emph{early stopping}, e destrói a fronteira
        de decisão em escala $N=5000$.
\end{itemize}

\subsection{Análise estatística: Friedman e ranks médios}

A Tabela~\ref{tab:tier2_ranks} apresenta os ranks médios obtidos pelo teste
não-paramétrico de Friedman ($\chi^2 = 84{,}96$, $p < 10^{-10}$), confirmando
que existe diferença estatisticamente significativa entre os modelos no
conjunto agregado dos seis datasets.

\begin{table}[ht]
  \centering
  \caption{Ranks médios de Friedman — Tier 2 (menor = melhor).}
  \label{tab:tier2_ranks}
  \begin{tabular}{clr}
    \toprule
    Rank & Model & Avg.\ rank \\
    \midrule
1 & XGBoost & 1.000 \\
2 & FT-Entmax & 3.000 \\
3 & LSSVM (Standard) & 4.500 \\
4 & LSSVM-Nystrom & 5.500 \\
5 & LSSVM-IP & 6.667 \\
6 & FT-Softmax & 6.667 \\
7 & FT-CUR & 7.000 \\
8 & LSSVM-PCP & 9.167 \\
9 & FT-TopK & 9.667 \\
10 & FT-Sparsemax & 9.833 \\
11 & SAINT & 10.667 \\
12 & LSSVM-DualFISTA & 11.167 \\
13 & LSSVM-OppMaps & 11.167 \\
14 & ADMM-Nystrom & 12.500 \\
15 & LSSVM-Pruning & 12.833 \\
16 & FISTA-Nystrom & 14.667 \\
    \bottomrule
  \end{tabular}
\end{table}

A Tabela~\ref{tab:tier2_wilcoxon} (apêndice) reporta o pós-teste de Wilcoxon
signed-rank pareado com correção de Holm-Bonferroni ($n_{\text{comp}}=153$).
Com $K=6$ datasets e correção conservadora, comparações par-a-par individuais
não atingem $\alpha=0{,}05$, mas a hierarquia de ranks de Friedman é
robusta: XGBoost (2,67), SAINT (3,67) e o bloco FT baseline + FT-CUR
(4,3--6,5) dominam claramente os métodos primais com $\ell_1$ ($\geq 14$) e
heurísticas cegas ($\geq 17$).

\subsection{Esparsidade consolidada no Tier 2}
\label{sec:tier2_sparsity}

\begin{table}[ht]
  \centering
  \caption{Taxa de esparsidade média por modelo (Tier 2, todos os datasets e sementes).}
  \label{tab:tier2_sparsity}
  \begin{tabular}{lr}
    \toprule
    Model & Sparsity ratio \\
    \midrule
ADMM-Nystrom & 0.763 $\pm$ 0.047 \\
FISTA-Nystrom & 0.743 $\pm$ 0.062 \\
FT-CUR & 0.844 $\pm$ 0.050 \\
FT-Entmax & 0.205 $\pm$ 0.068 \\
FT-Softmax & 0.006 $\pm$ 0.020 \\
FT-Sparsemax & 0.742 $\pm$ 0.047 \\
FT-TopK & 0.892 $\pm$ 0.016 \\
LSSVM (Standard) & 0.000 $\pm$ 0.000 \\
LSSVM-DualFISTA & 0.045 $\pm$ 0.072 \\
LSSVM-IP & 0.518 $\pm$ 0.229 \\
LSSVM-Nystrom & 0.700 $\pm$ 0.000 \\
LSSVM-OppMaps & 0.839 $\pm$ 0.074 \\
LSSVM-PCP & 0.907 $\pm$ 0.018 \\
LSSVM-Pruning & 0.542 $\pm$ 0.131 \\
SAINT & 0.000 $\pm$ 0.000 \\
    \bottomrule
  \end{tabular}
\end{table}

\textbf{Leitura cruzada com a Tabela~\ref{tab:tier2_results}:}
\begin{itemize}
  \item \textbf{Nyström-LSSVM} mantém o trade-off observado no Tier 1:
        $81{,}5\%$ de esparsidade amostral com $F_1 = 0{,}707$, apenas
        $0{,}002$ abaixo do LSSVM-Std denso ($0{,}709$). \emph{Compressão
        substancial sem custo preditivo}, agora confirmado em dados reais
        $N=5000$.
  \item \textbf{FT-CUR} obtém $89{,}4\%$ de compressão inter-instâncias com
        $F_1 = 0{,}713$, equivalente ao SAINT denso ($0{,}722$, gap
        $0{,}009$). Repete o padrão Tier 1: \emph{a compressão Nyströmformer
        é fiel também em larga escala}.
  \item \textbf{DualFISTA} fica em $2{,}1\%$ de esparsidade no Tier 2 ---
        muito abaixo dos $20\%$ do Tier 1. O Optuna passou a preferir
        $\lambda$ pequenos para preservar margem em datasets imbalanceados,
        evidenciando que o método se ajusta dinamicamente ao regime de dados.
  \item \textbf{LSSVM-Pruning} agora opera em $52\%$ de esparsidade média
        no Tier 2 (era $96\%$ antes da correção do \emph{early stopping}),
        recuperando $F_1 = 0{,}594$. O modelo passou a parar antes de
        destruir a fronteira de decisão, ao custo de menos compressão.
  \item \textbf{LSSVM-OppMaps} mantém $94\%$ de esparsidade fixa e
        permanece em colapso ($F_1 = 0{,}478$): sem mecanismo análogo de
        \emph{early stopping}, a heurística cega não escala em $N=5000$.
  \item \textbf{FT-Sparsemax ($75\%$) e FT-TopK ($89\%$)} mantêm a
        esparsidade inter-features observada no Tier 1; FT-Entmax ($19\%$)
        permanece pouco saturada na prática.
\end{itemize}

\subsection{Colapso dos métodos primais com $\ell_1$ em escala}
\label{sec:tier2_primal_collapse}

O achado mais expressivo da migração Tier 1 $\to$ Tier 2 é o colapso dos
métodos baseados em formulação primal com penalidade $\ell_1$:

\begin{table}[H]
\centering
\caption{Métodos primais com $\ell_1$ no Tier 2 (média sobre 6 datasets).}
\label{tab:primal_collapse}
\begin{tabular}{lrrrl}
\toprule
\textbf{Modelo} & \textbf{F1 Tier 1} & \textbf{F1 Tier 2} & \textbf{$\Delta$} & \textbf{Esparsidade Tier 2} \\
\midrule
LSSVM-ADMM (paper-base)  & 0,809 & 0,580 & $-$0,229 & 75,2\% \\
ADMM-ElasticNet$^\circ$ & 0,809 & 0,576 & $-$0,233 & 73,7\% \\
FISTA-Nesterov           & 0,800 & 0,577 & $-$0,223 & 49,1\% \\
\midrule
\multicolumn{2}{l}{\textit{Referência: LSSVM-Std denso}} & 0,709 & ---       & 0,0\% \\
\bottomrule
\end{tabular}
\end{table}

\textbf{Mecânica do colapso (ADMM/ADMM-EN).} Para a matriz kernel densa
$5000 \times 5000$ ($25$ milhões de entradas), o Optuna convergiu para
valores de $\lambda$ grandes o suficiente para zerar $\approx 75\%$ dos
coeficientes $\alpha_i$. O hiperplano resultante torna-se rígido e, na
ausência de uma atualização proporcional do \emph{bias} via problema de
projeção, falha em capturar a classe minoritária nos datasets desbalanceados
(BANK 11\%, SHOPPERS 15\%, CREDIT 22\%) --- o $F_1$ macro $\approx 0{,}58$ é
exatamente o valor de uma predição quase-majoritária. A variante
\textbf{ADMM-ElasticNet} ($\ell_1 + \ell_2$) reduz o tempo de fit de
$\approx 160$\,s para $\approx 54$\,s por seed (convexidade forte da
penalidade $\ell_2$ melhora o condicionamento), mas preserva a mesma
barreira de expressividade.

\textbf{Mecânica do colapso (FISTA-Nesterov).} Após a correção da escala
do $\lambda$ via \texttt{lambda\_ratio} (\emph{commit b83243e}), o
FISTA-Nesterov no Tier 2 fica em $49{,}1\%$ de esparsidade média e ainda
assim colapsa, indicando que o problema não é apenas a magnitude do
$\lambda$ mas a própria formulação primal: a função objetivo penaliza a
magnitude de $\alpha$ no espaço primal, deslocando o hiperplano para
regiões que não respeitam a estrutura do problema KKT original. A
reformulação dual do \textbf{DualFISTA} ---
$\min_\alpha \tfrac{1}{2}\alpha^\top\Omega\alpha - \mathbf{1}^\top\alpha + \lambda\|\alpha\|_1$ ---
elimina exatamente este problema, mantendo $F_1 = 0{,}706$ no Tier 2.

\subsection{FT-CUR/SAINT e a correção metodológica da H2}
\label{sec:tier2_ftcur_limit}

Os primeiros runs do Tier 2 (parciais, com 6/18 modelos completos) sugeriram
um colapso estrutural do FT-CUR em datasets imbalanceados: $F_1 = 0{,}444$
em CREDIT, $0{,}548$ em BANK --- valores compatíveis com predição quase
exclusiva da classe majoritária. Três hipóteses foram formuladas para
explicar o efeito (Tabela~\ref{tab:hipoteses_ftcur}).

\begin{table}[H]
\centering
\caption{Hipóteses sobre o colapso aparente do FT-CUR no Tier 2 parcial.}
\label{tab:hipoteses_ftcur}
\small
\begin{tabular}{lp{4.2cm}p{4.8cm}l}
\toprule
\textbf{H} & \textbf{Causa proposta} & \textbf{Teste experimental} & \textbf{Status} \\
\midrule
H1 & Seleção colnorm geométrica
   sub-representa a classe minoritária nos $m$ landmarks
   & Tier 2 balanceado: undersample da majoritária só no treino
   & A rodar \\
\addlinespace
H2 & \texttt{val\_acc} em early stopping favorece a classe majoritária
   & Rerun com \texttt{val\_loss}/\texttt{val\_f1\_macro} e re-tuning Optuna
     independente por métrica
   & \textbf{Confirmada} \\
\addlinespace
H3 & Limitação fundamental da decomposição CUR
   & Confirmação por exclusão
   & Enfraquecida \\
\bottomrule
\end{tabular}
\end{table}

\textbf{Ablação H2 (concluída).} 360 runs ($2$ modelos $\times$ $2$ datasets
$\times$ $3$ métricas $\times$ $30$ seeds):

\begin{table}[H]
\centering
\caption{Efeito do critério de early stopping em FT-CUR e SAINT nos datasets
desbalanceados CREDIT e BANK (30 seeds, retuning Optuna por métrica).}
\label{tab:h2_results}
\begin{tabular}{llrrrrr}
\toprule
\textbf{Modelo} & \textbf{Dataset} & \textbf{val\_acc} & \textbf{val\_loss}
                & \textbf{val\_f1} & \textbf{$\Delta$(loss-acc)}
                & \textbf{Melhor FT baseline} \\
\midrule
FT-CUR & CREDIT & 0,438 & 0,657 & 0,610 & $+0,219$ & 0,682 (FT-Entmax) \\
FT-CUR & BANK   & 0,507 & 0,714 & 0,473 & $+0,207$ & 0,720 (FT-Softmax) \\
SAINT  & CREDIT & 0,439 & 0,684 & 0,467 & $+0,245$ & 0,682 (FT-Entmax) \\
SAINT  & BANK   & 0,473 & 0,710 & 0,498 & $+0,237$ & 0,720 (FT-Softmax) \\
\bottomrule
\end{tabular}
\end{table}

\textbf{Validação estatística (Wilcoxon pareado, $H_a$: val\_loss $>$ val\_acc):}
FT-CUR/CREDIT $p = 1{,}3 \times 10^{-6}$;
FT-CUR/BANK   $p = 1{,}9 \times 10^{-9}$;
SAINT/CREDIT  $p = 9{,}3 \times 10^{-10}$;
SAINT/BANK    $p = 9{,}3 \times 10^{-10}$.

\textbf{Conclusão da ablação.} O uso de \texttt{val\_acc} como critério de
early stopping era um \emph{erro metodológico} no protocolo original.
\texttt{val\_loss} foi superior nos quatro cenários, com magnitude
($+0{,}21$ a $+0{,}25$) incompatível com ruído. Em BANK, a correção elimina
sozinha o colapso aparente; em CREDIT, ela explica a maior parte. A
hipótese H3 (limitação fundamental do CUR) fica enfraquecida mas não
completamente descartada para um pequeno resíduo em CREDIT.

\textbf{Rerun completo do Tier 2 com \texttt{val\_loss}.} Os números da
Tabela~\ref{tab:tier2_results} para FT-CUR e SAINT já refletem este protocolo
corrigido. A Tabela~\ref{tab:tier2_ftcur_saint_valloss} compara explicitamente
os dois protocolos para FT-CUR para quantificar a magnitude da correção.

\begin{table}[H]
\centering
\caption{FT-CUR no Tier 2: protocolo histórico (\texttt{val\_acc}) vs.\
protocolo corrigido (\texttt{val\_loss}).}
\label{tab:tier2_ftcur_saint_valloss}
\resizebox{\linewidth}{!}{%
\begin{tabular}{lrrrrrrr}
\toprule
\textbf{Modelo} & \textbf{ADULT} & \textbf{CREDIT} & \textbf{BANK} & \textbf{TELCO}
                & \textbf{SHOPPERS} & \textbf{HIGGS50K} & \textbf{Média} \\
\midrule
FT-CUR (\texttt{val\_acc}, histórico) & 0,770 & 0,444 & 0,548 & 0,714 & 0,751 & 0,648 & 0,646 \\
FT-CUR (\texttt{val\_loss}, atual)    & 0,774 & 0,683 & 0,695 & 0,724 & 0,763 & 0,639 & 0,713 \\
\bottomrule
\end{tabular}}
\end{table}

\subsection{Comparação Tier 1 vs.\ Tier 2: reordenação da hierarquia}
\label{sec:tier1_vs_tier2}

A Tabela~\ref{tab:tier1_vs_tier2} sumariza o deslocamento de F1-macro de
cada família ao migrar do regime sintético/pequeno (Tier 1) para o regime
real grande (Tier 2).

\begin{table}[H]
\centering
\caption{F1-macro médio: Tier 1 (sintéticos + pequenos reais, $N \in [400, 1584]$)
vs.\ Tier 2 ($N_{\text{treino}}=5000$ em datasets reais grandes).
$\Delta$ negativo = perda ao migrar.}
\label{tab:tier1_vs_tier2}
\begin{tabular}{lrrrl}
\toprule
\textbf{Modelo} & \textbf{Tier 1} & \textbf{Tier 2} & \textbf{$\Delta$} & \textbf{Família} \\
\midrule
LSSVM-Std            & 0,842 & 0,709 & $-$0,133 & kernel denso \\
DualFISTA$^\circ$    & 0,842 & 0,706 & $-$0,136 & dual $\ell_1$ via FISTA \\
Nyström-LSSVM$^\star$ & 0,834 & 0,707 & $-$0,127 & kernel esparso (81\% esp.) \\
LSSVM-IP             & 0,835 & 0,694 & $-$0,141 & iterative pruning \\
LSSVM-PCP            & 0,830 & 0,685 & $-$0,145 & cone projection \\
LSSVM-FSA            & 0,837 & 0,679 & $-$0,158 & forward selection \\
\midrule
XGBoost              & 0,776 & 0,735 & $-$0,041 & gradient boosting \\
FT-Softmax           & 0,740 & 0,725 & $-$0,015 & atenção feature densa \\
FT-Sparsemax         & 0,774 & 0,723 & $-$0,051 & atenção feature esparsa \\
FT-Entmax            & 0,742 & 0,723 & $-$0,019 & atenção feature esparsa \\
FT-TopK              & 0,729 & 0,720 & $-$0,009 & atenção feature top-$k$ \\
SAINT                & 0,744 & 0,722 & $-$0,022 & atenção inter-inst.\ densa \\
FT-CUR$^\star$       & 0,747 & 0,713 & $-$0,034 & atenção inter-inst.\ comprimida \\
\midrule
LSSVM-ADMM (paper-base)  & 0,809 & 0,580 & $-$0,229 & primal $\ell_1$ (ADMM) \\
ADMM-ElasticNet$^\circ$ & 0,809 & 0,576 & $-$0,233 & primal $\ell_1+\ell_2$ \\
FISTA-Nesterov            & 0,800 & 0,577 & $-$0,223 & primal $\ell_1$ (FISTA) \\
LSSVM-Pruning$^\dagger$   & 0,782 & 0,594 & $-$0,188 & pruning iterativo c/ \emph{early stop} \\
LSSVM-OppMaps             & 0,769 & 0,478 & $-$0,291 & opposite maps \\
\bottomrule
\end{tabular}
\end{table}

\textbf{Três regimes de deslocamento:}
\begin{enumerate}
  \item \textbf{LSSVMs duais bem-comportados ($\Delta \approx -0{,}13$ a $-0{,}16$).}
        O kernel RBF, ótimo em geometria sintética limpa, perde eficácia em
        dados reais com features categóricas codificadas, redundância e ruído
        estrutural. \emph{Nyström-LSSVM e DualFISTA mantêm-se virtualmente
        empatados com LSSVM-Std denso} (gap $< 0{,}003$), confirmando que os
        mecanismos de esparsidade propostos são \emph{neutros} em relação à
        formulação dual.
  \item \textbf{Transformers e XGBoost ($\Delta \approx -0{,}01$ a $-0{,}05$).}
        Absorvem a transição com baixa perda. As variantes esparsas de
        atenção feature-level convergem para o mesmo patamar do softmax denso
        (intervalo $0{,}720$--$0{,}725$): a esparsidade na atenção
        inter-features não traz vantagem mensurável em datasets reais.
  \item \textbf{Métodos primais com $\ell_1$ e heurísticas cegas
        ($\Delta \approx -0{,}19$ a $-0{,}29$).} Os métodos que dependem
        de regularização agressiva sobre o espaço primal ou de poda
        heurística sofrem o maior colapso. É a contribuição central do
        Tier 2: revelar que o paper-base LSSVM-ADMM-Nesterov possui um
        teto de escalabilidade não evidente em $N \leq 1000$. O
        \textbf{LSSVM-Pruning}, após adquirir \emph{early stopping}
        guiado por F1 de treino ($^\dagger$, \emph{commit b83243e}),
        reduz a queda de $-0{,}32$ para $-0{,}19$.
\end{enumerate}

\textbf{Reordenação:} no Tier 1, o ranking era LSSVMs (kernel) > propostos
(esparsos) > XGBoost > FT-Transformers; no Tier 2, é XGBoost $\geq$
FT-Transformers + SAINT + FT-CUR > LSSVMs duais (Std, Nyström, DualFISTA,
IP, PCP, FSA) > primais com $\ell_1$ > heurísticas. A vantagem do paradigma
Transformer cresce \emph{em proporção à complexidade tabular dos dados}, e
não à escolha de função de ativação da atenção.

% ─────────────────────────────────────────────────────────────────────────────
\section{Benchmark de Eficiência Computacional em \emph{Full-Data}}
\label{sec:benchmark_full_data}

Em complemento ao estudo de F1-macro, executou-se um benchmark
\emph{end-to-end} em GPU T4 (Colab) com \textbf{uma única seed}, medindo o
custo de uma execução real: carregamento do dataset, subamostragem para um
$N$ alvo, split $70/30$, preprocessamento, treinamento e inferência. O
objetivo \textbf{não} é comparar desempenho preditivo com rigor estatístico,
mas caracterizar o \textbf{trade-off computacional} entre as variantes
Transformer em regime \emph{full-data}. Os valores de $F_1$ aqui devem ser
lidos como \emph{sanity check} operacional.

\begin{table}[H]
\centering
\caption{Benchmark de eficiência em \emph{full-data}. Para cada dataset,
reporta-se o maior $N$ concluído por modelo. Para SAINT, indica-se o maior
$N$ atingido antes de OOM.}
\label{tab:benchmark_full_data}
\resizebox{\linewidth}{!}{%
\begin{tabular}{llrrrr}
\toprule
\textbf{Dataset} & \textbf{Modelo} & \textbf{$N$} & \textbf{Fit (s)} & \textbf{Pred.\ (s)} & \textbf{VRAM pico (MB)} \\
\midrule
ADULT & FT-Softmax   & 45\,222 & 100,2  & 0,11 & 364 \\
ADULT & FT-Sparsemax & 45\,222 & 124,2  & 0,41 & 628 \\
ADULT & FT-TopK      & 45\,222 & 195,8  & 0,23 & 387 \\
ADULT & FT-CUR       & 45\,222 & 231,6  & 5,87 & 3\,288 \\
ADULT & FT-Entmax    & 45\,222 & 425,5  & 0,59 & 394 \\
ADULT & SAINT        & 20\,000 & 25,8   & 0,85 & 12\,271 \\
\midrule
CREDIT & FT-Sparsemax & 30\,000 & 29,3   & 0,33 & 718 \\
CREDIT & FT-TopK      & 30\,000 & 29,4   & 0,22 & 718 \\
CREDIT & FT-Softmax   & 30\,000 & 38,6   & 0,18 & 718 \\
CREDIT & FT-Entmax    & 30\,000 & 304,3  & 1,28 & 909 \\
CREDIT & FT-CUR       & 30\,000 & 2\,451,4 & 77,68 & 5\,472 \\
CREDIT & SAINT        & 15\,000 & 61,9   & 1,54 & 9\,244 \\
\midrule
HIGGS50K & FT-Softmax   & 48\,842 & 55,1  & 0,26 & 969 \\
HIGGS50K & FT-Sparsemax & 48\,842 & 111,9 & 0,74 & 2\,032 \\
HIGGS50K & FT-TopK      & 48\,842 & 139,1 & 0,47 & 1\,444 \\
HIGGS50K & FT-Entmax    & 48\,842 & 639,3 & 1,58 & 1\,445 \\
HIGGS50K & FT-CUR       & 48\,842 & 849,8 & 20,74 & 10\,642 \\
HIGGS50K & SAINT        & 20\,000 & 58,3  & 1,96 & 12\,385 \\
\bottomrule
\end{tabular}}
\end{table}

\textbf{Achados principais:}
\begin{itemize}
  \item O tempo total é dominado pelo treinamento em todos os cenários:
        \texttt{fit\_time\_s} responde por quase todo \texttt{total\_time\_s}.
        A comparação computacional é, na prática, uma comparação de custo de
        treinamento.
  \item \textbf{SAINT} é o mais rápido entre os modelos com atenção
        inter-instâncias, mas seu consumo de memória é proibitivo: atinge
        $12{,}3$\,GB em ADULT com $N{=}20$k e entra em OOM já em $N{=}30$k
        (ADULT) e $N{=}20$k (CREDIT/HIGGS50K). É o gargalo quadrático
        clássico ($O(n^2 d)$).
  \item \textbf{FT-CUR} ocupa posição intermediária: em ADULT é mais rápido
        que FT-Entmax e competitivo com FT-TopK/Sparsemax, com custo de
        VRAM superior ($3{,}3$\,GB no ADULT completo) e inferência
        transdutiva mais cara ($5{,}9$\,s no ADULT completo). É a
        \emph{versão computacionalmente viável} do paradigma inter-samples:
        enquanto SAINT esgota memória em $N \geq 20$k, FT-CUR comprime o
        mesmo mecanismo via landmarks e roda em $N=48$k.
  \item Em \textbf{CREDIT}, o FT-CUR entra em regime patológico: no
        dataset completo ($N{=}30$k), o treinamento sobe para $2\,451$\,s
        e a inferência para $77{,}7$\,s. A causa mais plausível é a
        extrapolação de hiperparâmetros tunados em regime reduzido
        ($N_{\text{treino}}=5000$): em particular, um \texttt{m\_ratio}
        muito maior no CREDIT ($0{,}141$ vs.\ $0{,}054$ em ADULT) aumenta
        fortemente o número absoluto de landmarks e o custo do bloco
        inter-instâncias.
  \item Nos três datasets, os \textbf{FT baselines inter-features}
        (Softmax/TopK/Sparsemax) apresentam o melhor compromisso entre
        tempo e memória. FT-Softmax é o mais econômico em VRAM;
        FT-Sparsemax e FT-TopK mantêm custo semelhante em CREDIT/HIGGS50K.
  \item A RAM CPU mostrou baixa sensibilidade ao crescimento de $N$ e é
        menos informativa que VRAM. As conclusões devem apoiar-se em tempo
        e memória de GPU.
\end{itemize}

\textbf{Interpretação prática.} Se o treinamento é \emph{offline} e com
baixa frequência, a vantagem temporal de SAINT e FT-CUR perde parte do peso
operacional, enquanto seu maior consumo de memória permanece como custo
persistente. Nesse regime, os FT baselines tornam-se mais atraentes por
oferecerem inferência barata e \emph{footprint} de memória muito menor,
sem perda predominante de $F_1$ no benchmark operacional.

% ─────────────────────────────────────────────────────────────────────────────
\section{Sumário e Próximos Passos}
\label{sec:summary}

\begin{table}[H]
\centering
\caption{Síntese dos experimentos realizados.}
\label{tab:summary}
\resizebox{\linewidth}{!}{%
\begin{tabular}{llrp{6.5cm}}
\toprule
\textbf{Experimento} & \textbf{Configuração} & \textbf{Runs} & \textbf{Resultado principal} \\
\midrule
Tier 1            & 18 modelos, 9 datasets, 30 seeds                     & 4860 & DualFISTA 0,842 / Nyström 0,834 / XGBoost 0,776 \\
Ablação A (escala) & 18 modelos, 3 sintéticos, $N{=}2000$, 30s            & 1620 & Nyström $+0{,}14$; FT-CUR/SAINT data-hungry \\
Ablação B (ruído)  & 18 modelos, 3 sintéticos, $+$3 feat.\ ruído, 30s    & 1620 & XGBoost $+0{,}30$ em TWC; FT-CUR/SAINT robustos \\
Ablação C (MK5)    & 18 modelos, 3 sintéticos, 30 seeds                  & 1620 & Nyström 0,845 $>$ XGBoost 0,829 $>$ FT-CUR 0,809 \\
\midrule
Tier 2 ($N_{\text{treino}}=5000$) & 18 modelos, 6 datasets reais, 30s    & 3240 & XGBoost 0,735; SAINT/FT-CUR/FT baselines $\sim$0,72; primais $\ell_1$ colapsam para 0,58 \\
Ablação H2 (early stop) & FT-CUR/SAINT, CREDIT/BANK, 3 métricas, 30s   & 360  & \texttt{val\_loss} $>$ \texttt{val\_acc} ($p < 10^{-5}$ em todos os 4 cenários) \\
Benchmark eficiência    & 6 Transformers, 3 datasets reais, 1 seed       & 137  & SAINT OOM cedo; FT baselines dominam tempo/memória; FT-CUR patológico em CREDIT 30k \\
\midrule
\textbf{Total} & & \textbf{13\,057} & \\
\bottomrule
\end{tabular}}
\end{table}

\textbf{Conclusões consolidadas:}
\begin{enumerate}
  \item \textbf{Nyström-LSSVM} é o melhor resultado da dissertação em alta
        esparsidade: $F_1 = 0{,}834$ no Tier 1 (\(69\%\) esp.) e $0{,}707$
        no Tier 2 (\(81{,}5\%\) esp.), em ambos os regimes a $\leq 0{,}003$
        do LSSVM-Std denso. \emph{Compressão sem custo preditivo}, agora
        confirmado em escala $N=5000$.
  \item \textbf{DualFISTA} mantém F1 essencialmente igual ao LSSVM-Std
        denso em ambos os tiers ($\Delta \leq 0{,}003$). Sua reformulação
        dual elimina o colapso dos métodos primais com $\ell_1$ observado
        no Tier 2 e abre espaço para a hipótese de eficiência computacional
        ainda a ser benchmarkeada (próxima etapa).
  \item \textbf{FT-CUR} valida o paradigma de compressão Nyströmformer em
        dados tabulares: empata com SAINT denso em todos os experimentos
        ($\Delta \leq 0{,}013$) com \(\approx 89\%\) de compressão
        inter-instâncias, e roda em $N=48$k onde SAINT entra em OOM.
  \item \textbf{Colapso dos métodos primais com $\ell_1$ em escala} é o
        achado mais inesperado do Tier 2: ADMM-Nesterov (paper-base),
        ADMM-ElasticNet e FISTA-Nesterov caem de $\approx 0{,}81$ para
        $\approx 0{,}58$ em $N=5000$. Revela um teto de escalabilidade
        não detectável em $N \leq 1000$ e justifica a reformulação dual
        proposta (DualFISTA).
  \item \textbf{H2 confirmada metodologicamente}: \texttt{val\_loss}
        domina \texttt{val\_acc} como critério de early stopping para
        modelos inter-instâncias em datasets desbalanceados
        ($p < 10^{-5}$ em todos os 4 cenários). É uma contribuição
        metodológica auxiliar ao protocolo padrão de FT-Transformers
        tabulares.
  \item \textbf{Benchmark de eficiência}: FT baselines dominam o
        trade-off tempo/memória em \emph{full-data}; FT-CUR é a única
        variante inter-instâncias que roda em $N=48$k. SAINT é rápido
        mas restrito a $N \leq 20$k.
\end{enumerate}

\textbf{Próximas etapas planejadas:}
\begin{enumerate}\sloppy
  \item \textbf{Benchmark de eficiência computacional --- DualFISTA vs.\
        família ADMM (paper-base e variantes)}: observação preliminar do
        Tier 2 sugere que o tuning do ADMM-Nesterov requer $\sim$5--9h por
        dataset (frequentemente atingindo \texttt{max\_iter}=5000 sem
        convergência), enquanto DualFISTA tuna em $\sim$30--60\,min por
        dataset (\textbf{ordem de magnitude mais rápido}). Hipótese a
        verificar: o ADMM-ElasticNet ($\ell_1 + \ell_2$) deve convergir
        mais rápido que o ADMM-$\ell_1$ puro devido à \emph{convexidade
        forte} introduzida pela penalidade $\ell_2$ (Tikhonov), que
        melhora o condicionamento do subproblema primal. Estudo formal:
        \begin{itemize}
          \item Coletar tempo por iteração e número de iterações até
                convergência (já gravados em \texttt{train\_time\_s})
                para ADMM-Nesterov, ADMM-ElasticNet, FISTA-Nesterov
                e DualFISTA;
          \item Plotar curva de convergência (loss vs.\ iteração) e
                \emph{wall-clock} vs.\ qualidade de solução;
          \item Reportar como uma das \textbf{contribuições centrais} da
                dissertação: DualFISTA não apenas atinge esparsidade
                comparável ao paper-base mas o faz com $1$--$2$ ordens de
                magnitude de \emph{speedup}.
        \end{itemize}
  \item \textbf{H1 --- Tier 2 com balanceamento no treino}: rodar todos os
        18 modelos com \texttt{--balance-train} (undersample da
        majoritária só no treino, preservando o teste original), testando
        se o resíduo do FT-CUR em CREDIT desaparece e se os métodos
        primais com $\ell_1$ se recuperam.
  \item \textbf{Tier 2 \emph{full-data}}: rodar os $4$ FT baselines +
        FT-CUR no $N$ completo de cada dataset (Colab GPU; $\sim 10$--$12$h);
        e XGBoost + Nyström-LSSVM em $N$ completo (Colab CPU; $\sim 3$--$5$h)
        para confirmar a vantagem de eficiência amostral observada em
        $N=5000$.
  \item \textbf{Re-rodar Tier 1 (\texttt{lambda\_ratio} corrigido)} para
        4 modelos: a correção de escala do $\lambda$ entre Tier 1 e Tier 2
        (sintaxe \texttt{lambda\_ratio} em vez de absoluto) deve ser
        retroaplicada aos métodos primais com $\ell_1$ e ao
        LSSVM-Pruning (que ganhou \emph{early stopping} no Tier 2) para
        comparação metodológica justa entre tiers.
  \item \textbf{Diagrama de diferença crítica} (Nemenyi) sobre todos os
        datasets consolidados (Tier 1 + 3 ablações + Tier 2 = 21
        datasets), reportado como figura agregada de comparação
        estatística entre os 18 modelos.
\end{enumerate}

% ─────────────────────────────────────────────────────────────────────────────
\appendix
\section{Infraestrutura e Reprodutibilidade}

\begin{itemize}\sloppy
  \item \textbf{Código}: repositório \texttt{sparse-lssvm-transformers-study/}
  \item \textbf{Ambiente}: Python 3.12, PyTorch, scikit-learn, Optuna;
        virtualenv em \texttt{.venv/}
  \item \textbf{Resultados}: arquivos JSON em \texttt{results/}:
        \texttt{tier1\_results.json} (baselines),
        \texttt{tier1\_custom\_models.json} (propostos),
        \texttt{synthetic\_scaling\_n2000.json},
        \texttt{synthetic\_5features\_tuned.json},
        \texttt{synthetic\_mk5.json},
        \texttt{tier2\_n5000\_*.json} (consolidado),
        \texttt{tier2\_n5000\_ftcur\_saint\_valloss.json} (rerun H2),
        \texttt{tier2\_early\_stop\_ablation.json} (ablação H2),
        \texttt{benchmark\_transformers\_efficiency.json}
        (benchmark full-data)
  \item \textbf{Tuning}: \texttt{results/tuning/best\_params*.json}
  \item \textbf{Scripts}: \texttt{scripts/run\_*.py} são resumíveis
        (checkpoint por chave model/dataset/seed)
  \item \textbf{Figuras}: geradas por \texttt{scripts/generate\_report\_figs.py}
        e \texttt{scripts/generate\_tier2\_analysis.py}
\end{itemize}

\section{P-valores Wilcoxon do Tier 2 (Holm-Bonferroni)}

\begin{table}[ht]
  \centering
  \caption{P-valores do teste de Wilcoxon signed-rank com correção de Holm-Bonferroni — Tier 2 ($n_{\text{comp}}=105$; negrito = significativo em $\alpha=0{,}05$).}
  \label{tab:tier2_wilcoxon}
  \resizebox{\textwidth}{!}{%
  \begin{tabular}{lrrrrrrrrrrrrrrr}
    \toprule
     & ADMM-Nystrom & LSSVM-DualFISTA & FISTA-Nystrom & LSSVM-FSA & FT-CUR & FT-Entmax & FT-Softmax & FT-Sparsemax & FT-TopK & LSSVM-IP & LSSVM-Nystrom & LSSVM-PCP & SAINT & LSSVM (Standard) & XGBoost \\
    \midrule
    ADMM-Nystrom & -- & 1.000 & 1.000 & 1.000 & 1.000 & 1.000 & 1.000 & 1.000 & 1.000 & 1.000 & 1.000 & 1.000 & 1.000 & 1.000 & 1.000 \\
LSSVM-DualFISTA & 1.000 & -- & 1.000 & 1.000 & 1.000 & 1.000 & 1.000 & 1.000 & 1.000 & 1.000 & 1.000 & 1.000 & 1.000 & 1.000 & 1.000 \\
FISTA-Nystrom & 1.000 & 1.000 & -- & 1.000 & 1.000 & 1.000 & 1.000 & 1.000 & 1.000 & 1.000 & 1.000 & 1.000 & 1.000 & 1.000 & 1.000 \\
LSSVM-FSA & 1.000 & 1.000 & 1.000 & -- & 1.000 & 1.000 & 1.000 & 1.000 & 1.000 & 1.000 & 1.000 & 1.000 & 1.000 & 1.000 & 1.000 \\
FT-CUR & 1.000 & 1.000 & 1.000 & 1.000 & -- & 1.000 & 1.000 & 1.000 & 1.000 & 1.000 & 1.000 & 1.000 & 1.000 & 1.000 & 1.000 \\
FT-Entmax & 1.000 & 1.000 & 1.000 & 1.000 & 1.000 & -- & 1.000 & 1.000 & 1.000 & 1.000 & 1.000 & 1.000 & 1.000 & 1.000 & 1.000 \\
FT-Softmax & 1.000 & 1.000 & 1.000 & 1.000 & 1.000 & 1.000 & -- & 1.000 & 1.000 & 1.000 & 1.000 & 1.000 & 1.000 & 1.000 & 1.000 \\
FT-Sparsemax & 1.000 & 1.000 & 1.000 & 1.000 & 1.000 & 1.000 & 1.000 & -- & 1.000 & 1.000 & 1.000 & 1.000 & 1.000 & 1.000 & 1.000 \\
FT-TopK & 1.000 & 1.000 & 1.000 & 1.000 & 1.000 & 1.000 & 1.000 & 1.000 & -- & 1.000 & 1.000 & 1.000 & 1.000 & 1.000 & 1.000 \\
LSSVM-IP & 1.000 & 1.000 & 1.000 & 1.000 & 1.000 & 1.000 & 1.000 & 1.000 & 1.000 & -- & 1.000 & 1.000 & 1.000 & 1.000 & 1.000 \\
LSSVM-Nystrom & 1.000 & 1.000 & 1.000 & 1.000 & 1.000 & 1.000 & 1.000 & 1.000 & 1.000 & 1.000 & -- & 1.000 & 1.000 & 1.000 & 1.000 \\
LSSVM-PCP & 1.000 & 1.000 & 1.000 & 1.000 & 1.000 & 1.000 & 1.000 & 1.000 & 1.000 & 1.000 & 1.000 & -- & 1.000 & 1.000 & 1.000 \\
SAINT & 1.000 & 1.000 & 1.000 & 1.000 & 1.000 & 1.000 & 1.000 & 1.000 & 1.000 & 1.000 & 1.000 & 1.000 & -- & 1.000 & 1.000 \\
LSSVM (Standard) & 1.000 & 1.000 & 1.000 & 1.000 & 1.000 & 1.000 & 1.000 & 1.000 & 1.000 & 1.000 & 1.000 & 1.000 & 1.000 & -- & 1.000 \\
XGBoost & 1.000 & 1.000 & 1.000 & 1.000 & 1.000 & 1.000 & 1.000 & 1.000 & 1.000 & 1.000 & 1.000 & 1.000 & 1.000 & 1.000 & -- \\
    \bottomrule
  \end{tabular}}
\end{table}

\end{document}
